% concatenated from FAST.tex
\documentclass[10pt, reqno]{amsart}
%\documentclass[10pt, reqno, draft]{amsart}
%\documentclass[a4paper, 10pt]{tran-l}
\usepackage{epsfig, amssymb, amsmath, version}
\usepackage{amssymb, version, graphicx, fancybox, mathrsfs, multirow}
\usepackage{epstopdf}
\usepackage{url, hyperref}
% \usepackage[notcite, notref]{showkeys}
\usepackage{bm}
\usepackage{subcaption}
\usepackage{float}

\usepackage{color, xcolor}
\usepackage{cases}
\usepackage{mathtools}
\usepackage{extarrows}
\usepackage{bbm}
\DeclarePairedDelimiter{\ceil}{\lceil}{\rceil}

\textheight=21.6cm
\textwidth=14.2cm
%\textwidth=18cm
\setlength{\oddsidemargin}{1cm}
\setlength{\evensidemargin}{1cm}
%---- Basic definition ----
\catcode`\@=11 \theoremstyle{plain}
\@addtoreset{equation}{section}   % Makes \section reset 'equation' counter. 
\renewcommand{\theequation}{\arabic{section}.\arabic{equation}}
\@addtoreset{figure}{section}
\renewcommand\thefigure{\thesection. \@arabic\c@figure}
\renewcommand{\thefigure}{\arabic{section}.\arabic{figure}}
%\@addtoreset{theorem}{section}
\newtheorem{thm}{\bf Theorem}
\renewcommand{\thethm}{\arabic{section}.\arabic{thm}}


\newenvironment{theorem}{\begin{thm}} {\end{thm}}
\newtheorem{cor}{\bf Corollary}
\newtheorem{prop}{Proposition}[section]
\renewcommand{\thecor}{\arabic{section}.\arabic{cor}}
\newenvironment{corollary}{\begin{cor}} {\end{cor}}
\newtheorem{lmm}{\bf Lemma}
\renewcommand{\thelmm}{\arabic{section}.\arabic{lmm}}
\newenvironment{lemma}{\begin{lmm}}{\end{lmm}}
\newcommand{\sign}{\text{sign}}
\DeclareMathOperator{\sgn}{sign}

%\theoremstyle{example}
%\newtheorem{example}{Example}[section]


\theoremstyle{remark}
\newtheorem{rem}{\bf Remark}[section]
%\theoremstyle{conjecture}
\newtheorem{conjecture}{\bf Conjecture}[section]
\theoremstyle{definition}
\newtheorem{defn}{\bf Definition}[section]
\newtheorem{prob}{\bf Problem}[section]
\renewcommand{\thetable}{\thesection.\arabic{table}}
\numberwithin{table}{section}



\newcommand{\re}[1]{(\ref{#1})}
\newcommand{\ree}[2]{(\ref{#1})--(\ref{#2})}
\def \ri {{\rm i}}
\def \af {\alpha}
\def \bt {\beta}

\allowdisplaybreaks[4]
%\renewcommand  \varphi \phi

\newcommand{\bs}[1]{\boldsymbol{#1}}
\def \er {{\bs e}_r}
\def \et {{\bs e}_\theta}
\def \ep {{\bs e}_\varphi}
\def \vt {{\bs T}_l^m}
\def \vv {{\bs V}_l^m}
\def \vw {{\bs W}_l^m}
\def \rd {{\rm d}}
\def \R {{\mathbb R}}

\renewcommand \wedge \times

\newcommand{\comm}[1]{\marginpar{%
		\vskip-\baselineskip %raise the marginpar a bit
		\raggedright\footnotesize
		\itshape\hrule\smallskip#1\par\smallskip\hrule}}


\newcommand{\diff}[1]{\ensuremath{\mathrm{d}#1}}
\newcommand{\Question}[1]{%
	\marginpar{\flushleft\scriptsize\bfseries\upshape#1}}

\def \dlx {D_{-}^{s}}

\def \drx {D_{+}^{s}}

\newcommand {\cred}[1] {{\color{red} #1}}

\begin{document}
	\bibliographystyle{plain}
	\graphicspath{{./figures/}}
	%\baselineskip 13pt
	
\title[Fast Finite Difference Scheme for Tempered Time Fractional]
	 {A New Fast Finite Difference Scheme for Tempered Time Fractional Advection-Dispersion Equation with a Weak Singularity at Initial Time}
	\author[Liangcai Huang]
{\;\; Liangcai Huang${}^1$,  \;\; Shujuan Lü${}^2$ }
	
	\thanks{${}^{1}$School of Mathematics,  Beihang University(BUAA),  Beijing,  102206,  China.  Email: hlcsy@buaa.edu.cn (L. Huang).  \\
	 \indent ${}^2$Corresponding author.  School of Mathematics,  Beihang University(BUAA),  Beijing,  102206,  China.  Email: lsj@buaa.edu.cn (S. Lü). 
			}

\begin{abstract}
In this paper,  we propose a new second-order fast finite difference scheme in time for solving the Tempered Time Fractional Advection-Dispersion Equation.  Under the assumption that the solution is nonsmooth at the initial time,  we investigate the uniqueness,  stability,  and convergence of the scheme.  Furthermore,  we prove that the scheme achieves second-order convergence in both time and space.  Finally,  corresponding numerical examples are provided. 
\end{abstract}
\keywords{Caputo-tempered fractional derivative $\cdot$ Weak Singularity $\cdot$ Variable Time Stepping $\cdot$ Fast Finite difference method $\cdot$ Stability $\cdot$ Convergence}
%
% \subjclass[2010]{65N35,  65E05,  65N12,   41A10,  41A25,  41A30,  41A58}	

\maketitle


%\vspace*{-12pt}


\section{Introduction}
Tempered fractional calculus is an advanced extension of classical fractional calculus,  designed to model complex phenomena characterized by anomalous diffusion and nonlocal behavior.  It introduces a tempered waiting time distribution to address scenarios where classical models may fail.  Koponen (1995)\cite{koponen1995} proposed the truncated Lévy flight with a smooth cutoff as an optimization of Mantegna's (1994)\cite{mantegna1994} truncated Lévy flight,  offering improved properties.  Subsequently,  Cartea and del Castillo-Negrete (2007)\cite{cartea2007} developed the tempered fractional calculus operator based on fractional calculus to address the infinite moment problem associated with power-law waiting times,  ensuring the finiteness of moments in the Lévy distribution.  To date,  tempered fractional calculus has been widely applied in various fields,  including finance\cite{carr2002, carr2003, zhou2023} ,  geophysics\cite{meerschaert2008, xia2013, zhang2011, zhang2012} and groundwater hydrology\cite{zhang2014}.  By capturing the complexity of memory effects and non-Gaussian dynamics,  tempered fractional calculus enhances our understanding and predictive capabilities of complex systems. 

The application of tempered fractional calculus to practical problems requires numerical algorithms capable of solving tempered fractional differential equations.  Currently,  tempered fractional calculus has been applied to various models,  including the Bloch equation\cite{feng2022},  the tempered fractional Burgers equation\cite{dwivedi2025, wang2022}, the tempered fractional Jacobi function\cite{zhao2023} and the temporal tempered fractional Feynman-Kac equation\cite{zhao2024}.  Among these,  the numerical solution of Caputo Time-Tempered Fractional Diffusion Equations (TTFDE) has been extensively studied in several articles\cite{fenwick2024, dwivedi2024, krzyzanowski2024, morgado2019, safari2022, zhao2020, zhou2023}.  The TTFDE is defined as follows
\begin{align}
{ }_0^C D_t^{\alpha,  \lambda} u=\frac{\partial^2u}{\partial x^2}-\frac{\partial u}{\partial x}+f(x,  t),  \quad(x,  t) \in \Omega \times(0,  T] \label{no1}
\end{align}
where $\lambda > 0,  0 < \alpha < 1, \Omega=[0, L]$. 

For $0<\alpha<1$,  the Caputo tempered fractional derivative can be denoted as
$$
{ }_0^C D_t^{\alpha,  \lambda} u(t)=\frac{e^{-\lambda t}}{\Gamma(1-\alpha)} \int_0^t \frac{1}{(t-s)^{\alpha}}\left(e^{\lambda s} u(s)\right)^{'} d s
$$
where $\Gamma(\cdot)$ means the Gamma function. 
And when $\lambda = 0$, the Caputo tempered fractional derivative is equal to the Caputo fractional derivative which is defined as
$$
{ }_0^C D_t^\alpha u(t)=\frac{1}{\Gamma(1-\alpha)} \int_0^t \frac{1}{(t-s)^{\alpha}} u^{'}(s) d s. 
$$

In the study of TTFDE,  various numerical methods have been developed to solve tempered fractional differential equations,  achieving different levels of accuracy depending on the method and grid type: Morgado (2019)\cite{morgado2019} utilized the finite difference method to solve tempered fractional advection-diffusion equations.  By employing a variable grid for nonsmooth solutions,  the method achieved a temporal accuracy of $\min\{2 - \alpha,  r\alpha\}$ and a spatial accuracy of second order. Zhao (2020)\cite{zhao2020} applied the finite difference method to numerically solve the time fractional sub-diffusion equation.  Using a uniform grid for smooth solutions,  the method achieved a temporal accuracy of $2 - \alpha$ and a spatial accuracy of second order. Safari (2022)\cite{safari2022} used the local discontinuous Galerkin (LDG) method to solve time–space tempered fractional diffusion equations.  By adopting a nonuniform grid for nonsmooth solutions,  the method achieved a temporal accuracy of $\min(2 - \alpha,  r\delta)$ and high-order spatial accuracy. Zhou (2023)\cite{zhou2023} proposed a fast finite difference method for solving tempered time-fractional diffusion-reaction equations.  With a variable grid for nonsmooth solutions,  the method achieved a temporal accuracy of $\min\{2 - \alpha,  r\alpha\}$ and a spatial accuracy of fourth order. Krzyżanowski (2024)\cite{krzyzanowski2024} applied the finite difference method to solve the fractional tempered diffusion problem.  Using a uniform grid for smooth solutions,  the method achieved a temporal accuracy of $2 - \alpha$ and a spatial accuracy of second order. Fenwick (2024)\cite{fenwick2024} developed a fast finite difference method to solve the fractional tempered diffusion problem.  With a uniform grid for smooth solutions,  the method achieved a temporal accuracy of $2 - \alpha$ and a spatial accuracy of second order. Dwivedi (2024)\cite{dwivedi2024} utilized a fast finite difference method to solve the 2D tempered fractional reaction-advection-subdiffusion equation.  Using a uniform grid for smooth solutions,  the method achieved a temporal accuracy of $1 + \alpha$ and a spatial accuracy of fourth order. These studies demonstrate the effectiveness of various numerical approaches in achieving high accuracy for solving tempered fractional differential equations under different conditions. 

In this paper,  inspired by the Tempered Time
Fractional Advection-Dispersion Equation(TTFADE) presented by Xia\cite{xia2013} and Meerschaert\cite{meerschaert2008}. We develop a high-order difference schemes for the tempered Time Fractional Advection-Dispersion Equation
\begin{align}
&\frac{\partial u}{\partial t}+{ }_0^C D_t^{\alpha,  \lambda} u=\frac{\partial^2u}{\partial x^2}-\frac{\partial u}{\partial x}+f(x,  t),  \quad(x,  t) \in \Omega \times(0,  T]\label{Equation1}\\
& u(x,  0)=\phi(x),  \quad x \in \Omega \label{Equation2}\\
&u(x,  t)|_{X \in \partial \Omega}=0,  \quad t \in[0,  T]\label{Equation3}
\end{align}
where $\lambda > 0,  0 < \alpha < 1, \Omega=[0, L]$. 

Unlike TTFDE,  the TTFADE studied in this paper additionally includes a first-order time derivative term, it is highly significant for assessing and restoring the groundwater environment. In existing literature,  Cao(2020)\cite{cao2020} employed a fast finite difference-finite element method to solve the TTFADE, using a variable grid to handle nonsmooth solutions,  the method achieved first-order accuracy in time and high-order accuracy in space. Qiao(2022)\cite{qiao2022} utilized a second-order backward difference formula and a second-order convolution quadrature method for time discretization to solve the multi-dimensional tempered fractional integro-differential equation.  By employing a uniform grid to solve smooth solutions,  the method achieved a temporal accuracy of \(1 + \alpha\) and a spatial accuracy of second order. 

As it is well know, the challenging part of numerically solving the equation arises from the weakly singular kernel and exponential kernel involved in the tempered fractional operator. This study aims to contribute to this evolving field by proposing a novel numerical scheme that integrates fast finite difference methods with temporal 2-order schemes on adaptive grids for solving the TTFADE under the assumption that the solution is nonsmooth at the initial time. We will analyze the stability and convergence of our method,  highlight its advantages over existing approaches,  and validate its effectiveness through a series of numerical experiments.  

\section{Fast Evaluation of the Tempered Caputo fractional Derivative}
In this section,  we give a quick calculation format for the tempered Caputo fractional derivation with $0\leq \alpha \leq 1$,  assuming that the time fractional derivative is defined on [0,  T],  using a gradient grid $T_N=\{t_n|t_n=T(\frac{n}{N})^r, n = 0,  1,  \cdots ,  N \}$, where N is a positive integer,  $r \geq 1$ and the subinterval length $\tau_n = t_n-t_{n-1}$. Let $t_{n+\frac{1}{2}}=\frac{1}{2}(t_n+t_{n+1})$. 

\subsection{Derivation of Fast Approaching Operator}
\begin{lemma}\cite{beylkin2005}\label{fast}[\textbf{Fast approximation of power function}]\label{fast}
For any $\alpha \geq 0$ , the absolute error $\varepsilon $ and  $\delta= \tau_1 $,  there are positive integer $N_{exp} $, the set of positive numbers $\{s_l | l = 1,  2,  \cdots,  N_ {exp} \} $ and the corresponding positive weights $\{\omega_l | l = 1, 2, \cdots, N_{exp}\}$ such that
$$|t^{-1-\alpha}-\sum_{l=1}^{N_{exp}}\omega_le^{-s_lt}|\leqslant\varepsilon\quad t\in[\delta, T], $$
where
$$N_{exp}=\mathcal{O}((log\frac1\varepsilon)(loglog\frac1\varepsilon+log\frac  T\delta)+(log\frac1\delta)(loglog\frac1\varepsilon+log\frac1\delta))$$
and the positive real numbers $s_i$,  $w_i, (i = 1,  2,  \cdots ,  N_{exp})$ and $N_{exp}$ can be obtained from \cite{beylkin2005}. 
\end{lemma}

Firstly,  a discrete approximation to the fractional derivative ${ }_0^C D_t^{\alpha,  \lambda} u(t)$ at $t_{n+\frac{1}{2}}$ can be divided into a history part $C_h\left(t_n\right)$ and a local part $C_l\left(t_n\right)$ by the following approximation
\begin{equation}
\begin{aligned}
\left. { }_0^C D_t^{\alpha,  \lambda} u(t)\right|_{t=t_{n+\frac{1}{2}}} & =\left. e^{-\lambda t_{n+\frac{1}{2}}} {}_0^CD_t^\alpha\left(e^{\lambda t} u(t)\right)\right|_{t=t_{n+\frac{1}{2}}} \\
& =\frac{e^{-\lambda t_{n+\frac{1}{2}}}}{\Gamma(1-\alpha)} \int_0^{t_{n+\frac{1}{2}}} \frac{1}{\left(t_{n+\frac{1}{2}}-s\right)^\alpha}\left(e^{\lambda s} u(s)\right)^{\prime} d s \\
& =\frac{e^{-\lambda t_{n+\frac{1}{2}}}}{\Gamma(1-\alpha)}\left[\int_0^{t_{n}} \frac{1}{\left(t_{n+\frac{1}{2}}-s\right)^\alpha} d\left(e^{\lambda s} u(s)\right)+\int_{t_{n}}^{t_{n+\frac{1}{2}}} \frac{1}{\left(t_{n+\frac{1}{2}}-s\right)^\alpha} d\left(e^{\lambda s} u(s)\right)\right] \\
&= \frac{e^{-\lambda t_{n+\frac{1}{2}}}}{\Gamma(1-\alpha)}\int_0^{t_{n}} \frac{1}{\left(t_{n+\frac{1}{2}}-s\right)^\alpha} d\left(e^{\lambda s} u(s)\right)+\frac{e^{-\lambda t_{n+\frac{1}{2}}}}{\Gamma(1-\alpha)}\int_{t_{n}}^{t_{n+\frac{1}{2}}} \frac{1}{\left(t_{n+\frac{1}{2}}-s\right)^\alpha} d\left(e^{\lambda s} u(s)\right)\\
& \triangleq{} C_h\left(t_{n+\frac{1}{2}}\right)+C_l\left(t_{n+\frac{1}{2}}\right) . 
\end{aligned}
\end{equation}

For the history part,  via applying Lemma \ref{fast} to replace the convolution kernel $\frac{1}{t^{1+\alpha}}$,  it follows from the equatioan above that
\begin{equation}
\begin{aligned}\label{Ch}
C_h\left(t_{n+\frac{1}{2}}\right) & =\frac{e^{-\lambda t_{n+\frac{1}{2}}}}{\Gamma(1-\alpha)} \int_0^{t_{n}} \frac{1}{\left(t_{n+\frac{1}{2}}-s\right)^\alpha} d\left(e^{\lambda s} u(s)\right) \\
& =\frac{e^{-\lambda t_{n+\frac{1}{2}}}}{\Gamma(1-\alpha)}\left[\frac{e^{\lambda s} u(s)}{\left(t_{n+\frac{1}{2}}-s\right)^\alpha} |_{t_0}^{t_{n}}-\alpha \int_0^{t_{n}} e^{\lambda s} u(s) \frac{1}{\left(t_{n+\frac{1}{2}}-s\right)^{\alpha+1}} d s\right] \\
& =\frac{e^{-\lambda t_{n+\frac{1}{2}}}}{\Gamma(1-\alpha)}\left[\frac{e^{\lambda t_{n}} u\left(t_{n}\right)}{(\frac{\tau_{n+1}}{2})^\alpha}-\frac{u\left(t_0\right)}{t_{n+\frac{1}{2}}^\alpha}-\alpha \int_0^{t_{n}} \frac{e^{\lambda s} u(s)}{\left(t_{n+\frac{1}{2}}-s\right)^{\alpha+1}} d s\right] \\
& \approx \frac{e^{-\lambda t_{n+\frac{1}{2}}}}{\Gamma(1-\alpha)}\left[\frac{e^{\lambda t_{n}} u\left(t_{n}\right)}{(\frac{\tau_{n+1}}{2})^\alpha}-\frac{u\left(t_0\right)}{t_{n+\frac{1}{2}}^\alpha}-\alpha \sum_{i=1}^{N_{\text {exp }}} \omega_i e^{\lambda t_{n+\frac{1}{2}}} U_{h i s,  i}\left(t_n\right)\right] \\
& =\frac{1}{\Gamma(1-\alpha)}\left[\frac{e^{-\lambda \frac{\tau_{n+1}}{2}} u\left(t_{n}\right)}{(\frac{\tau_{n+1}}{2})^\alpha}-\frac{e^{-\lambda t_{n+\frac{1}{2}}}}{t_{n+\frac{1}{2}}^\alpha} u\left(t_0\right)-\alpha \sum_{i=1}^{N_{\text {exp }}} \omega_i U_{h i s,  i}\left(t_n\right)\right]
\end{aligned}
\end{equation}

with $U_{h i s,  i}\left(t_n\right)=\int_0^{t_{n}} e^{-\left(\lambda+s_i\right)\left(t_{n+\frac{1}{2}}-s\right)} u(s) d s$.  By the definition of $U_{h i s,  i}\left(t_n\right)$,  it is not difficult to observe that
$$
U_{h i s,  i}\left(t_n\right)=e^{-\left(\lambda+s_i\right) \frac{\tau_n+\tau_{n+1}}{2}} U_{h i s,  i}\left(t_{n-1}\right)+\int_{t_{n-1}}^{t_{n}} e^{-\left(\lambda+s_i\right)\left(t_{n+\frac{1}{2}}-s\right)} u(s) d s . 
$$
Replacing the exact solution function $u$ by a linear interpolation function $L_{1,  n} u \triangleq{} \frac{s-t_{n-1}}{\tau_{n}}u(t_{n})+\frac{t_{n}-s}{\tau_{n}}u(t_{n-1})$ on the subinterval $\left[t_{n-1},  t_{n}\right]$,  let 
\begin{equation}\label{lambda}
    \begin{aligned}
        \lambda_{i, n}^1&=\int_{t_{n-1}}^{t_{n}} e^{-\left(\lambda+s_i\right)\left(t_{n+\frac{1}{2}}-s\right)} \frac{s-t_{n-1}}{\tau_{n}} d s\\
        &=\frac{e^{-(\lambda+s_i)\frac{\tau_{n+1}}{2}}}{(\lambda+s_i)^2\tau_{n}}\left(e^{-\left(\lambda+s_i\right) \tau_{n}}-1+\left(\lambda+s_i\right) \tau_{n}\right), \\
        \lambda_{i, n}^{2}&=\int_{t_{n-1}}^{t_{n}} e^{-\left(\lambda+s_i\right)\left(t_{n+\frac{1}{2}}-s\right)} \frac{t_{n}-s}{\tau_{n}} d s\\
&=\frac{e^{-(\lambda+s_i)\frac{\tau_{n+1}}{2}}}{(\lambda+s_i)^2\tau_{n}}\left(1-e^{-\left(\lambda+s_i\right) \tau_{n}}-e^{-\left(\lambda+s_i\right) \tau_{n}}\left(\lambda+s_i\right) \tau_{n}\right). 
    \end{aligned}
\end{equation}
we could compute the integral
\begin{equation}
\begin{aligned}
\int_{t_{n-1}}^{t_{n}} e^{-\left(\lambda+s_i\right)\left(t_{n+\frac{1}{2}}-s\right)} u(s) d s \approx \lambda_{i, n}^{1}u(t_{n})+\lambda_{i, n}^{2}u(t_{n-1}). 
\end{aligned}
\end{equation}
So 
$$
\begin{aligned}
    U_{his, i}(t_n) &\approx e^{-\left(\lambda+s_i\right) \frac{\tau_n+\tau_{n+1}}{2}} U_{h i s,  i}\left(t_{n-1}\right)+\lambda_{i, n}^{1}u(t_{n})+\lambda_{i, n}^{2}u(t_{n-1})\\
    &=\sum_{j=0}^{n-1} e^{-\left(\lambda+s_i\right)\left(t_{n+\frac{1}{2}}-t_{{n+\frac{1}{2}}-j}\right)}\left(\lambda_{i,  n-j}^1 u(t_{n-j})+\lambda_{i,  n-j}^2 u(t_{n-1-j})\right). 
\end{aligned}
$$
Then we can simplify representation of the term $\alpha \sum_{i=1}^{N_{\text {exp }}} \omega_i U_{h i s,  i}\left(t_n\right)$ in (\ref{Ch}), such that
\begin{equation}
\begin{aligned}
\alpha &\sum_{i=1}^{N_{\text {exp }}} \omega_i U_{h i s,  i}\left(t_n\right) \\
& =\sum_{j=0}^{n-1} \alpha \sum_{i=1}^{N_{e x p}} \omega_i e^{-\left(\lambda+s_i\right)\left(t_{n+\frac{1}{2}}-t_{{n+\frac{1}{2}}-j}\right)}\left(\lambda_{i,  n-j}^1 u(t_{n-j})+\lambda_{i,  n-j}^2 u(t_{n-1-j})\right) \\
& =\sum_{j=0}^{n-1}\left(a_{j,  n} u(t_{n-j})+b_{j,  n} u(t_{n-1-j})\right)\\
&=\sum_{l=1}^{n-1}(a_{n-l, n}+b_{n-1-l, n})u(t_l)+a_{0, n}u(t_{n})+b_{n-1, n}u(t_0). 
\end{aligned}
\end{equation}
where
\begin{equation}\label{ab}
\begin{aligned}
&a_{j,  n}=\alpha \sum_{i=1}^{N_{\exp }} \omega_i e^{-\left(\lambda+s_i\right)\left(t_{n+\frac{1}{2}}-t_{{n+\frac{1}{2}}-j}\right)} \lambda_{i,  n-j}^1,  \\
&b_{j,  n}=\alpha \sum_{i=1}^{N_{\exp }} \omega_i e^{-\left(\lambda+s_i\right)\left(t_{n+\frac{1}{2}}-t_{{n+\frac{1}{2}}-j}\right)} \lambda_{i,  n-j}^2. 
\end{aligned}
\end{equation}
Then we get
$$
\begin{aligned}
    C_h(t_{n+\frac{1}{2}}) \approx \frac{1}{\Gamma(1-\alpha)}\left[\frac{e^{-\lambda \frac{\tau_{n+1}}{2}} u\left(t_{n}\right)}{(\frac{\tau_{n+1}}{2})^\alpha}-\frac{e^{-\lambda t_{n+\frac{1}{2}}}}{t_{n+\frac{1}{2}}^\alpha} u\left(t_0\right)-\right. \\
    \left. \sum_{l=1}^{n-1}(a_{n-l, n}+b_{n-1-l, n})u(t_l)-a_{0, n}u(t_{n})-b_{n-1, n}u(t_0)\right]. 
\end{aligned}
$$

As for the local part,  by an approach similar to the $L 1$ algorithm in \cite{sun2006},  the function $u$ is approximated by its linear interpolation function $L_{1,  n+1} u$,  namely, 
\begin{equation}
\begin{aligned} \label{Cl}
C_l\left(t_{n+\frac{1}{2}}\right) & \approx \frac{e^{-\lambda t_{n+\frac{1}{2}}}}{\Gamma(1-\alpha)} \int_{t_{n}}^{t_{n+\frac{1}{2}}} \frac{1}{\left(t_{n+\frac{1}{2}}-s\right)^\alpha}\left(e^{\lambda s} L_{1,  n+1} u(s)\right)^{\prime} d s \\
& \approx \frac{e^{-\lambda t_{n+\frac{1}{2}}}}{\Gamma(1-\alpha)} \frac{e^{\lambda t_{n+\frac{1}{2}}} L_{1,  n+1} u\left(t_{n+\frac{1}{2}}\right)-e^{\lambda t_{n}} L_{1,  n+1} u\left(t_{n}\right)}{\frac{\tau_{n+1}}{2}} \int_{t_{n}}^{t_{n+\frac{1}{2}}} \frac{1}{\left(t_{n+\frac{1}{2}}-s\right)^\alpha} d s \\
&=\frac{e^{-\lambda t_{n+\frac{1}{2}}}}{\Gamma(1-\alpha)} \frac{e^{\lambda t_{n+\frac{1}{2}}}\frac{u(t_{n})+u(t_{n+1})}{2}-e^{\lambda t_{n}}u\left(t_{n}\right)}{\frac{\tau_{n+1}}{2}}\frac{(\frac{\tau_{n+1}}{2})^{1-\alpha}}{1-\alpha}\\
& =\frac{1}{\Gamma(1-\alpha)}\left[\frac{u\left(t_{n}\right)+u\left(t_{n+1}\right)}{2(1-\alpha)(\frac{\tau_{n+1}}{2})^\alpha} -\frac{e^{-\lambda \frac{\tau_{n+1}}{2}}}{(1-\alpha)(\frac{\tau_{n+1}}{2})^\alpha} u\left(t_{n}\right)\right]. 
\end{aligned}
\end{equation}

Synthesize the above formula, then we get
\begin{equation}
\begin{aligned}
{}^C_0D^{\alpha, \lambda}_tu(t)|_{t=t_{n+\frac{1}{2}}}\approx &\frac{1}{\Gamma(1-\alpha)}\left[\frac{u\left(t_{n}\right)+u\left(t_{n+1}\right)}{2(1-\alpha)(\frac{\tau_{n+1}}{2})^\alpha}-(a_{0, n}+\frac{\alpha e^{-\lambda \frac{\tau_{n+1}}{2}} }{(1-\alpha)(\frac{\tau_{n+1}}{2})^\alpha})u\left(t_{n}\right)-\right. \\
    &\left. \sum_{l=1}^{n-1}(a_{n-l, n}+b_{n-1-l, n})u(t_l)-(b_{n-1, n}+\frac{e^{-\lambda t_{n+\frac{1}{2}}}}{t_{n+\frac{1}{2}}^\alpha})u(t_0)\right]\\
    \triangleq& { }_0^{FC}D_t^{\alpha,  \lambda} u^{n+\frac{1}{2}}
\end{aligned}
\end{equation}

\subsection{Error analysis}
Define mesh functions
$$
u^n=u(t_n), 1\leq n \leq N. $$
From section 2. 1, our scheme of evaluating the Tempered Caputo fractional derivative is defined by ${ }_0^{FC}D_t^{\alpha,  \lambda} u^{n+\frac{1}{2}}$. 

The following three lemmas are the result of error estimates,  and they play an important role in analyzing the stability and convergence of our numerical methods.  In order to analyze the approximation error of fast evalution,  the L1-approximation of Caputo tempered fractional derivative $_0^C\mathbb{D}_t^{\alpha, \lambda} u^n$ is introduced firstly. 
\begin{equation}
\begin{aligned}
    _0^C\mathbb{D}_t^{\alpha, \lambda} u^{n+\frac{1}{2}}=&\frac{e^{-\lambda t_{n+\frac{1}{2}}}}{\Gamma(1-\alpha)}\sum_{k=0}^{n-1}\int_{t_k}^{t_{k+1}}\frac1{(t_{n+\frac{1}{2}}-s)^\alpha}\left(e^{\lambda s}L_{1, k+1}u(s)\right)^{\prime}ds\\
    &+\frac{e^{-\lambda t_{n+\frac{1}{2}}}}{\Gamma(1-\alpha)}\int_{t_n}^{t_{n+\frac{1}{2}}}\frac1{(t_{n+\frac{1}{2}}-s)^\alpha}\left(e^{\lambda s}L_{1, n+1}u(s)\right)^{\prime}ds
\end{aligned}
\end{equation}
where $L_{1, k+1}u$ is the linear interpolation function of $u(t)$ on the subinterval $\left[t_k, t_{k+1}\right]$. 
\begin{lemma}[\cite{cao2020}]\label{tau}
With the grading parameter $r \geq 1$,  the smoothly graded mesh is $t_k=T(k/N)^r$.  The following rewarding results hold and $\tau_k$ is monotonically increasing with respect to k,  i.e. 
$$\tau_{k+1}=T\left(\frac{k+1}N\right)^r-T\left(\frac kN\right)^r\leq C TN^{-r}k^{r-1}$$
$$\tau_{k+1}\geq \tau_k$$
for $k = 1, \cdots, N-1$. 
\end{lemma}
\begin{lemma}[\cite{diethelm2010}D. 6]\label{Gamma}Let $1-\alpha, \delta-1 \in R_+$. Then
    $$
\int_0^x (x-s)^{-\alpha}x^{\delta-2} ds=x^{\delta-\alpha-1} \frac{\Gamma(1-\alpha) \Gamma(\delta-1)}{\Gamma(\delta-\alpha)}
$$
\end{lemma}
\begin{lemma} \label{Ru}
Suppose that u(t) satisfies $\left|\frac{\partial^l u(t)}{\partial t^l}\right| \leq C(1+t^{\delta-l}),  1<\delta<2,  l=0, 1, 2$ and $0< \alpha < 1$.  Then there exists a constant K such that
\begin{equation}
    \begin{aligned}
        \left| R^{n+\frac{1}{2}}u\right|&\triangleq{}\left|{}^C_0D^{\alpha, \lambda}_tu(t)|_{t=t_{n+\frac{1}{2}}}- {}^C_0\mathbb{D}_t^{\alpha,  \lambda} u^{n+\frac{1}{2}}\right|\\
        &\leq K(n+1)^{-min\{2-\alpha, r(1+\delta-\alpha)\}},  n\geq 1. \\
|R^{\frac{1}{2}}u|&\leq KN^{-r(\delta+1-\alpha)}. \label{L1half}
\end{aligned}
\end{equation}
\end{lemma}
\begin{proof}
For any $n \geq 0$, 
    \begin{align}\label{R1}
        R^{n+\frac{1}{2}}u=&{}^C_0D^{\alpha, \lambda}_tu(t)|_{t=t_{n+\frac{1}{2}}}- {}^C_0\mathbb{D}_t^{\alpha,  \lambda} u^{n+\frac{1}{2}}\notag\\
        =&\frac{e^{-\lambda t_{n+\frac{1}{2}}}}{\Gamma(1-\alpha)}\sum_{k=0}^{n-1}\int_{t_k}^{t_{k+1}}\frac1{(t_{n+\frac{1}{2}}-s)^\alpha}\left(e^{\lambda s}u(s) -e^{\lambda s}L_{1, k+1}u(s)\right)^{\prime}ds\notag\\
        &+\frac{e^{-\lambda t_{n+\frac{1}{2}}}}{\Gamma(1-\alpha)}\int_{t_n}^{t_{n+\frac{1}{2}}}\frac1{(t_{n+\frac{1}{2}}-s)^\alpha}\left(e^{\lambda s}u(s)-e^{\lambda s}L_{1, n+1}u(s)\right)^{\prime}ds\notag\\
        =&\frac{\alpha e^{-\lambda t_{n+\frac{1}{2}}}}{\Gamma(1-\alpha)}\sum_{k=0}^{n-1}\int_{t_k}^{t_{k+1}}\frac{e^{\lambda s}u(s)-e^{\lambda s}L_{1, k+1}u(s)}{(t_{n+\frac{1}{2}}-s)^{\alpha+1}}ds\notag \\
        &+\frac{e^{-\lambda t_{n+\frac{1}{2}}}}{\Gamma(1-\alpha)}\int_{t_n}^{t_{n+\frac{1}{2}}}\frac{\lambda e^{\lambda s}u(s)+e^{\lambda s}u'(s)-\lambda e^{\lambda s}L_{1, n+1}u(s)-e^{\lambda s}L_{1, n+1}'u(s)}{(t_{n+\frac{1}{2}}-s)^\alpha}ds\notag\\
        \triangleq{}&K\sum_{k=0}^{n-1}\int_{t_k}^{t_{k+1}}\frac{e^{\lambda s-\lambda t_{n+\frac{1}{2}}}(u(s)-L_{1, k+1}u(s))}{(t_{n+\frac{1}{2}}-s)^{\alpha+1}}ds\notag\\
        &+\frac{1}{\Gamma(1-\alpha)}\int_{t_n}^{t_{n+\frac{1}{2}}}\frac{\lambda e^{\lambda s-\lambda t_{n+\frac{1}{2}}}[u(s)-L_{1, n+1}u(s)]+e^{\lambda s-\lambda t_{n+\frac{1}{2}}}[u'(s)-L_{1, n+1}'u(s)]}{(t_{n+\frac{1}{2}}-s)^\alpha}ds\notag\\
        \leq&K(\sum_{k=0}^{n-1}R_{{n+\frac{1}{2}}, k}+R_{{n+\frac{1}{2}}, n})
    \end{align}
where $K = \frac{1}{\Gamma(1-\alpha)}$ is a constant. 
When $k=0$, 
\begin{equation}
    \begin{aligned}
|R_{{n+\frac{1}{2}}, 0}| &= |\int_{t_0}^{t_{1}}\frac{e^{\lambda( s- t_{n+\frac{1}{2}})}(u(s)-L_{1, 1}u(s))}{(t_{n+\frac{1}{2}}-s)^{\alpha+1}}ds|\\
&\leq \int_{t_0}^{t_{1}}\frac{|u(s)-L_{1, 1}u(s)|}{(t_{n+\frac{1}{2}}-s)^{\alpha+1}}ds\\
& =\int_{t_0}^{t_1}\left(t_{n+\frac{1}{2}}-s\right)^{-(\alpha+1)}\left|\int_0^s\left(u(\theta)-L_{1, 1} u(\theta)\right)^{\prime} d \theta\right| d s \\
& \leq \int_{t_0}^{t_1}\left(t_{n+\frac{1}{2}}-s\right)^{-(\alpha+1)}\left(\int_0^s\left|u^{\prime}(\theta)\right| d \theta+\left|\int_0^s t_1^{-1}(u(t_1)-u(t_0)) d \theta\right|\right) d s \\
& \leq C\left(t_{n+\frac{1}{2}}-t_1\right)^{-(\alpha+1)}\left(\int_{t_0}^{t_1} s^\delta+s t_1^{\delta-1} d s\right) \\
& \leq C\left(t_{n+\frac{1}{2}}-t_1\right)^{-(\alpha+1)} (\frac{t_1^{\delta+1}}{\delta+1}+\frac{t_1^{\delta+1}}{2})\\
&\leq K\left(t_{n+\frac{1}{2}}-t_1\right)^{-(\alpha+1)}t_1^{\delta+1}\\
& \leq K\left(\frac{t_{n+\frac{1}{2}}-t_1}{t_1}\right)^{-(\alpha+1)}t_1^{\delta-\alpha} \\
& \leq K(n+1)^{-r(\alpha+1)}t_1^{\delta-\alpha}\\
&\triangleq{}K(n+1)^{-r(\alpha+1)}N^{-r(\delta-\alpha)}\\
&\leq K(n+1)^{-r(\delta+1)}
\end{aligned}
\end{equation}
In virtue of the fact that
\begin{align}
    \left| u(s)-L_{1,  k+1} u(s)\right|&=\frac{1}{2} \left|u^{\prime \prime}\left(\xi_{k+1}\right)\left(s-t_{k}\right)\left(t_{k+1}-s\right)\right|\notag\\
    &\leq \frac{1}{2}\max_{t \in\left[t_{k},  t_{k+1}\right]}|u''(t)|\left(\frac{t_{k+1}-t_k}{2}\right)^2\notag\\
    &\leq \frac{\tau_{k+1}^2}{8}Ct_k^{\delta-2}\notag\\
    &\triangleq{}K_1\tau_{k+1}^2t_k^{\delta-2} \label{e1}
\end{align}
for $1\leq k\leq \lceil \frac{n+1}{2}\rceil-1$, Lemma \ref{tau} give
    \begin{align*}
    |R_{n+\frac{1}{2}, k}| &= |\int_{t_k}^{t_{k+1}}\frac{e^{\lambda( s- t_{n+\frac{1}{2}})}(u(s)-L_{1, k+1}u(s))}{(t_{n+\frac{1}{2}}-s)^{\alpha+1}}ds|\\
    &\leq \int_{t_k}^{t_{k+1}}|\frac{u(s)-L_{1, k+1}u(s)}{(t_{n+\frac{1}{2}}-s)^{\alpha+1}}|ds \\
    &\leq K_1\tau_{k+1}^2t_k^{\delta-2}\int_{t_k}^{t_{k+1}}|\frac{1}{(t_{n+\frac{1}{2}}-s)^{\alpha+1}}|ds\\
    &\leq K_1\tau_{k+1}^3t_k^{\delta-2}\frac{1}{(t_{n+\frac{1}{2}}-t_{k+1})^{\alpha+1}}\\
    &\leq K_2(TN^{-r}k^{r-1})^3[T(\frac{k}{N})^r]^{\delta-2}[\frac{T(\frac{n+1}{N})^r+T(\frac{n}{N})^r}{2}-T(\frac{k+1}{N})^r]^{-(\alpha+1)}\\
    &=K_2T^{\delta-\alpha}N^{r(\alpha-\delta)}k^{r(1+\delta)-3}[\frac{(n+1)^r+n^r}{2}-(k+1)^r]^{-(1+\alpha)}\\
    &\leq K_2T^{\delta-\alpha}(n+1)^{r(\alpha-\delta)}k^{r(1+\delta)-3}[\frac{(n+1)^r+n^r}{2}-(\lceil \frac{n+1}{2}\rceil)^r]^{-(1+\alpha)}\\
    &\leq K_3(n+1)^{r(\alpha-\delta)}k^{r(1+\delta)-3}[(n+1)^r]^{-(1+\alpha)}\\
    &=K_3(n+1)^{r(\alpha-\delta)}k^{r(1+\delta)-3}\\
    &\triangleq{}K_3k^{r(1+\delta)-3}(n+1)^{-r(1+\delta)}
    \end{align*}
where $K_4$ is a constant. As in \cite{gracia2018}, we can use the inequality
\begin{equation}
    (n+1)^{-r(\delta+1)}\sum_{k=1}^{\lceil \frac{n+1}{2}\rceil-1}k^{r(\delta+1)-3}\leq C\begin{cases}(n+1)^{-r(\delta+1)}&\text{ if }r(\delta+1)<2, \\(n+1)^{-2}\ln n&\text{ if }r(\delta+1)=2, \\(n+1)^{-2}&\text{ if }r(\delta+1)>2, \end{cases}
\end{equation}
to bound the local sum in (\ref{R1}). 
For $\lceil \frac{n+1}{2}\rceil\leq k \leq n-1$, invoke Lemma \ref{tau}, and $t_{n+1}\geq t_k\geq2^{-r}t_{n+1}$, 
\begin{align*}
\left|\sum_{k=\lceil \frac{n+1}{2}\rceil}^{n-1}R_{{n+\frac{1}{2}}, k}\right|& \leq \sum_{k=\lceil \frac{n+1}{2}\rceil}^{n-1}K_1\tau_{k+1}^2t_k^{\delta-2}\int_{t_k}^{t_{k+1}}|\frac{1}{(t_{n+\frac{1}{2}}-s)^{\alpha+1}}|ds\\
    &\leq \sum_{k=\lceil \frac{n+1}{2}\rceil}^{n-1}\alpha K_1\tau_{n}^2(2^{-r}t_{n+1})^{\delta-2}[(t_{n+\frac{1}{2}}-t_{k+1})^{-\alpha}-(t_{n+\frac{1}{2}}-t_k)^{-\alpha}]\\
    &\triangleq{} K_2\tau_{n}^2t_{n+1}^{\delta-2}\sum_{k=\lceil n/2\rceil}^{n-1}[(t_{n+\frac{1}{2}}-t_{k+1})^{-\alpha}-(t_{n+\frac{1}{2}}-t_k)^{-\alpha}]\\
    &\leq K_2\tau_{n}^2t_{n+1}^{\delta-2}(t_{n+\frac{1}{2}}-t_{n})^{-\alpha}\\
    &= K_2 \tau_{n}^{2}t_{n+1}^{\delta-2}(\frac{\tau_{n+1}}{2})^{-\alpha}\\
    &\leq K_3 \tau_{n+1}^{2-\alpha}t_{n+1}^{\delta-2}\\
    &\leq  K_3(TN^{-r}(n+1)^{r-1})^{2-\alpha}[T(\frac{n+1}{N})^r]^{\delta-2}\\
    &\leq KN^{r(\alpha-\delta)}(n+1)^{r(\delta-\alpha)+\alpha-2}\\
    &\leq K(n+1)^{-(2-\alpha)}
\end{align*}
For $R_{{n+\frac{1}{2}}, n}(n\geq 1)$, 
\begin{align}
u(t_{n+1})&=u(s)+u'(s)(t_{n+1}-s)+\int_s^{t_{n+1}}u''(t)(t_{n+1}-t)dt\notag\\
u(t_{n})&=u(s)+u'(s)(t_{n}-s)+\int_s^{t_{n}}u''(t)(t_{n}-t)dt\notag\\
    \left|u'(s)-L_{1, n+1}'u(s)\right|&=\left|u'(s)-\frac{u(t_{n+1})-u(t_n)}{\tau_{n+1}}\right|\notag\\
    &=\left|\frac{\int_s^{t_{n+1}}u''(t)(t_{n+1}-t)dt-\int_s^{t_{n}}u''(t)(t_{n}-t)dt}{\tau_{n+1}}\right|\notag\\
    &=\left|[\int_s^{x}u''(t)(x-t)dt]'|_{x=\xi}\right|\notag\\
    &=\left|[\int_s^{x}u''(t)dt]|_{x=\xi}\right|\notag\\
    &=\left|\int_s^{\xi}u''(t)dt\right|, \xi\in [t_n, t_{n+1}] \notag\\
    &\leq \tau_{n+1}\max_{t\in[t_n, t_{n+1}]}|u''(t)|\notag\\
    &\leq C\tau_{n+1}t_n^{\delta-2}\label{e2}
\end{align}
Applying the inequalities (\ref{e1}) and (\ref{e2}), we have
    \begin{align}
\left|R_{{n+\frac{1}{2}}, n}\right|&\leq \int_{t_n}^{t_{n+\frac{1}{2}}}\frac{\lambda e^{\lambda s-\lambda t_{n+\frac{1}{2}}}\left|u(s)-L_{1, n+1}u(s)\right|+e^{\lambda s-\lambda t_{n+\frac{1}{2}}}\left|u'(s)-L_{1, n+1}'u(s)\right|}{(t_{n+\frac{1}{2}}-s)^\alpha}ds\notag\\
&\leq \int_{t_n}^{t_{n+\frac{1}{2}}}\frac{\lambda \left|u(s)-L_{1, n+1}u(s)\right|+\left|u'(s)-L_{1, n+1}'u(s)\right|}{(t_{n+\frac{1}{2}}-s)^\alpha}ds\notag\\
&\leq \int_{t_n}^{t_{n+\frac{1}{2}}}\frac{\lambda K_1\tau_{n+1}^2t_{n}^{\delta-2}+C\tau_{n+1}t_n^{\delta-2}}{(t_{n+\frac{1}{2}}-s)^\alpha}ds\notag\\
&\triangleq{}K_2\tau_{n+1}^{3-\alpha}t_{n}^{\delta-2}+K_3\tau_{n+1}^{2-\alpha}t_{n}^{\delta-2}\notag\\
&\leq K_2\left[TN^{-r}n^{r-1}\right]^{3-\alpha}\left[T\left(\frac{n}N\right)^r\right]^{\delta-2}+ K_3\left[TN^{-r}n^{r-1}\right]^{2-\alpha}\left[T\left(\frac{n}N\right)^r\right]^{\delta-2} \notag\\
&= K_2 T^{\delta-\alpha+1}N^{-r(\delta-\alpha+1)}\left(n-1\right)^{r(\delta-\alpha+1)-(3-\alpha)}+K_3 T^{\delta-\alpha}N^{-r(\delta-\alpha)}\left(n-1\right)^{r(\delta-\alpha)-(2-\alpha)} \notag\\
&\leq K_2T^{\delta-\alpha+1}n^{-(3-\alpha)}+K_3T^{\delta-\alpha}n^{-(2-\alpha)}\notag\\
&\leq K (n+1)^{-(2-\alpha)}. 
\end{align}
And for $R_{{n+\frac{1}{2}}, n}(n=0)$, we estimate two equations first with $s \in [t_0, t_\frac{1}{2}]$, 
\begin{align}
    |u(s)-L_{1, 1}u(s)|=&\left| u(s)-[\frac{t_1-s}{\tau_1}u(t_0)+\frac{s-t_0}{\tau_1}u(t_1)]\right|\notag\\
    =&\left|\frac{t_1-s}{\tau_1}[u(s)-u(t_0)]+\frac{s-t_0}{\tau_1}[u(s)-u(t_1)]\right|\notag\\
    =&\left|\frac{(t_1-s)(s-t_0)}{\tau_1}u'(\xi_1)+\frac{(s-t_0)(s-t_1)}{\tau_1}u'(\xi_2)\right|\notag\\
    \leq&2\tau_1\max_{t\in[t_0, t_1]}|u'(t)|\notag\\
    \leq &2\tau_1t_1^{\delta-1}\notag\\
    =&2t_1^{\delta}\label{1e}\\
    u'(s)-\frac{t(t_1)-u(t_0)}{\tau_1}=&\frac{1}{\tau_1}\left(\int_s^{t_{1}}u''(t)(t_{1}-t)dt-\int_s^{t_{0}}u''(t)(t_{0}-t)dt\right)\label{2e}
\end{align}
Applying the formula (\ref{1e}) and (\ref{2e}) , we get
\begin{align}
|R_{{\frac{1}{2}}, 0}|=& \left|\int_{t_0}^{t_{\frac{1}{2}}}\frac{(e^{\lambda s}u(s)-e^{\lambda s}L_{1, 1}u(s))'}{(t_{\frac{1}{2}}-s)^\alpha}ds\right|\notag\\
=&\left|\int_{t_0}^{t_{\frac{1}{2}}}\frac{\lambda e^{\lambda s}(u(s)-L_{1, 1}u(s))+e^{\lambda s}(u'(s)-\frac{t(t_1)-u(t_0)}{\tau_1})}{(t_{\frac{1}{2}}-s)^\alpha}ds\right| \notag\\
\leq &\left|\int_{t_0}^{t_{\frac{1}{2}}}\frac{\lambda e^{\lambda s}(u(s)-L_{1, 1}u(s))}{(t_{\frac{1}{2}}-s)^\alpha}ds\right| \notag\\
&+\left|\int_{t_0}^{t_{\frac{1}{2}}}\frac{e^{\lambda s}}{\tau_1}\left(\int_s^{t_{1}}u''(t)(t_{1}-t)dt-\int_s^{t_{0}}u''(t)(t_{0}-t)dt\right)\frac{ds}{(t_{\frac{1}{2}}-s)^\alpha}\right|\notag\\
\leq &\lambda e^{\lambda t_{\frac{1}{2}}}\int_{t_0}^{t_{\frac{1}{2}}}\frac{2t_1^\delta}{(t_{\frac{1}{2}}-s)^\alpha}ds \notag\\
&+\left|\int_{t_0}^{t_{\frac{1}{2}}}\frac{e^{\lambda s}}{\tau_1}\left(t_1\int_s^{t_{1}}u''(t)dt-\int_{t_0}^{t_{1}}tu''(t)dt\right)\frac{ds}{(t_{\frac{1}{2}}-s)^\alpha}\right|\notag\\
\leq& Kt_1^{\delta+1-\alpha}+\frac{e^{\lambda t_\frac{1}{2}}}{\tau_1}\int_{t_0}^{t_{\frac{1}{2}}}\left(t_1\int_s^{t_{1}}|u''(t)|dt+\int_{t_0}^{t_{1}}t|u''(t)|dt\right)\frac{ds}{(t_{\frac{1}{2}}-s)^\alpha}\notag\\
\leq &Kt_1^{\delta+1-\alpha}+\frac{e^{\lambda t_\frac{1}{2}}}{\tau_1}\int_{t_0}^{t_{\frac{1}{2}}}\left(t_1\int_s^{t_{1}}Ct^{\delta-2}dt+\int_{t_0}^{t_{1}}tCt^{\delta-2}dt\right)\frac{ds}{(t_{\frac{1}{2}}-s)^\alpha}\notag\\
\leq&Kt_1^{\delta+1-\alpha}+\frac{e^{\lambda t_\frac{1}{2}}}{\tau_1}\int_{t_0}^{t_{\frac{1}{2}}}\left(t_1Ct_1s^{\delta-2}+\frac{Ct_1^\delta}{\delta}\right)\frac{ds}{(t_{\frac{1}{2}}-s)^\alpha}\notag\\
\leq &Kt_1^{\delta+1-\alpha}+K_1\int_{t_0}^{t_{\frac{1}{2}}}t_1^2s^{\delta-2}(t_{\frac{1}{2}}-s)^{-\alpha}ds+K_2t_1^{\delta+1-\alpha}\notag\\
\triangleq{}&Kt_1^{\delta+1-\alpha}+Kt_1^2 t_1^{\delta-\alpha-1}\label{gamma}\\
\leq &KN^{-r(\delta+1-\alpha)}
\end{align}
where (\ref{gamma}) is obtained by Lemma \ref{Gamma}. 
Collecting the bounds above, we get 
\begin{equation}
        |R^{n+\frac{1}{2}}u| =|K(\sum_{k=0}^{n-1}R_{{n+\frac{1}{2}}, k}+R_{{n+\frac{1}{2}}, n})|\leq K(n+1)^{-min\{2-\alpha, r(1+\delta-\alpha)\}}, 
\end{equation}
which is the desired result. 
\end{proof}
Note that (11) is a modification of a result in \cite{cao2023},  where error is estimated as $n^{-min\{2-\alpha, r(1+\alpha)\}}$. 

\begin{lemma} \label{FRu}
Suppose that u(t) satisfies $\left|\frac{\partial^l u(t)}{\partial t^l}\right| \leq C(1+t^{\delta-l}),  1<\delta<2,  l=0, 1, 2$ and $0< \alpha < 1$. Then there exists a constant K which is independent of n such that
\begin{equation}
    \begin{aligned}
        |{}^{FC}R^{n+\frac{1}{2}}u|&\triangleq{}|{}^C_0D^{\alpha, \lambda}_tu(t)|_{t=t_{n+\frac{1}{2}}}- { }_0^{FC}D_t^{\alpha,  \lambda} u^{n+\frac{1}{2}}|\\
        &\leq K[(n+1)^{-min\{2-\alpha, r(1+\delta-\alpha)\}}+\varepsilon],  n\geq1. \\
        |{}^{FC}R^{\frac{1}{2}}u|&\leq K[N^{-r(\delta+1-\alpha)}+\varepsilon].  \label{Fhalf}
    \end{aligned}
\end{equation}
\end{lemma}
\begin{proof}
    For $n\geq1$, applying \ref{Ch} and \ref{Cl} we have
    \begin{equation}
    \begin{aligned}
        {}^{FC}R^{n+\frac{1}{2}}u&={}^C_0D^{\alpha, \lambda}_tu(t)|_{t=t_{n+\frac{1}{2}}}- {}^C_0\mathbb{D}_t^{\alpha,  \lambda} u^{n+\frac{1}{2}}+{}^C_0\mathbb{D}_t^{\alpha,  \lambda} u^{n+\frac{1}{2}}-{ }_0^{FC}D_t^{\alpha,  \lambda} u^{n+\frac{1}{2}}\\
        &=R^{n+\frac{1}{2}}u+{}^C_0\mathbb{D}_t^{\alpha,  \lambda} u^{n+\frac{1}{2}}-{ }_0^{FC}D_t^{\alpha,  \lambda} u^{n+\frac{1}{2}}\\
    \end{aligned}
\end{equation}
\begin{align*}
    \left|{}^C_0\mathbb{D}_t^{\alpha,  \lambda} u^{n+\frac{1}{2}}\right. &\left. - { }_0^{FC}D_t^{\alpha,  \lambda}u^{n+\frac{1}{2}}\right|\\
    \leq& \frac{e^{-\lambda t_{n+\frac{1}{2}}}}{\Gamma(1-\alpha)}\left[\sum_{k=0}^{n-1}\alpha\int_{t_k}^{t_{k+1}}e^{\lambda s}\left|L_{1, {k+1}}u(s)(\frac{1}{\left(t_{n+\frac{1}{2}}-s\right)^{\alpha+1}}-\sum_{l=1}^{N_{exp}}w_le^{-s_l(t_{n+\frac{1}{2}}-s)})\right|ds\right]\\
         &+\frac{e^{-\lambda t_{n+\frac{1}{2}}}}{\Gamma(1-\alpha)} \int_{t_{n}}^{t_{n+\frac{1}{2}}} \left\lvert\,  \frac{e^{\lambda t_{n+\frac{1}{2}}} L_{1,  n+1} u\left(t_{n+\frac{1}{2}}\right)-e^{\lambda t_n} L_{1,  n+1} u\left(t_n\right)}{\frac{\tau_{n+1}}{2}}\right.  \\
& \left. -\left(e^{\lambda s} L_{1,  n+1} u(s)\right)^{\prime}\right| \left(t_{n+\frac{1}{2}}-s\right)^{-\alpha} d s\\
\leq &\frac{e^{-\lambda t_{n+\frac{1}{2}}} \alpha \varepsilon}{\Gamma(1-\alpha)} \sum_{k=0}^{n-1} \int_{t_{k}}^{t_{k+1}} e^{\lambda s}\left|L_{1,  k+1} u(s)\right| d s+\\
&+\frac{e^{-\lambda t_{n+\frac{1}{2}}}}{\Gamma(1-\alpha)} \int_{t_{n}}^{t_{n+\frac{1}{2}}} \left\lvert\,  \frac{e^{\lambda t_{n+\frac{1}{2}}} L_{1,  n+1} u\left(t_{n+\frac{1}{2}}\right)-e^{\lambda t_n} L_{1,  n+1} u\left(t_n\right)}{\frac{\tau_{n+1}}{2}}\right.  \\
& \left. -\left(e^{\lambda s} L_{1,  n+1} u(s)\right)^{\prime}\right| \left(t_{n+\frac{1}{2}}-s\right)^{-\alpha} d s\\
\triangleq{}& {}^{FC}R^{n+\frac{1}{2}}_hu+ {}^{FC}R^{n+\frac{1}{2}}_lu
\end{align*}
Firstly, for ${}^{FC}R^{n+\frac{1}{2}}_hu$, applying the Lemma \ref{fast}
\begin{align}
        \left|{}^{FC}R^{n+\frac{1}{2}}_hu\right|&\leq \frac{ e^{-\lambda t_{n+\frac{1}{2}}}\alpha \varepsilon}{\Gamma(1-\alpha)}\sum_{k=0}^{n-1}\int_{t_k}^{t_{k+1}}e^{\lambda s}\left|L_{1, k+1}u(s) \right|ds \notag\\
&\leq \frac{\alpha e^{-\lambda t_{n+\frac{1}{2}}}\varepsilon}{\Gamma(1-\alpha)}\max_{s\in [0, t_n]}|u(s)|\int_0^{t_n}e^{\lambda s}ds\notag\\
&\leq \frac{\alpha \varepsilon}{\Gamma(1-\alpha)} t_n^\delta\int_0^{t_n}1ds \notag\\
&\leq \frac{\alpha \varepsilon}{\Gamma(1-\alpha)}t_n^{\delta+1}\notag\\
&\leq \frac{\alpha \varepsilon}{\Gamma(1-\alpha)}T^{\delta+1}\notag\\
&\triangleq{} K\varepsilon
\end{align}
Secondly, we can use the Taylor expansion to $(e^{\lambda s}L_{1, n+1}u(s))'$ and get
\begin{align*}
e^{\lambda t_{n+\frac{1}{2}}} L_{1,  n+1} u\left(t_{n+\frac{1}{2}}\right)=&e^{\lambda s}L_{1, n+1}u(s)+(e^{\lambda s}L_{1, n+1}u(s))'(t_{n+\frac{1}{2}}-s)\\
&+\int_{s}^{t_{n+\frac{1}{2}}}(e^{\lambda \theta}L_{1, n+1}u(\theta))''(t_{n+\frac{1}{2}}-\theta)d\theta
\end{align*}
\begin{align*}
e^{\lambda t_{n}} L_{1,  n+1} u\left(t_{n}\right)=&e^{\lambda s}L_{1, n+1}u(s)+(e^{\lambda s}L_{1, n+1}u(s))'(t_{n}-s)\\
&+\int_{s}^{t_{n}}(e^{\lambda \theta}L_{1, n+1}u(\theta))''(t_{n}-\theta)d\theta
\end{align*}
\begin{equation}
    \begin{aligned}\label{fe1}
        &\left|\frac{e^{\lambda t_{n+\frac{1}{2}}} L_{1,  n+1} u\left(t_{n+\frac{1}{2}}\right)-e^{\lambda t_{n}} L_{1,  n+1} u\left(t_{n}\right)}{\frac{\tau_{n+1}}{2}} -(e^{\lambda s}L_{1, n+1}u(s))'\right|\\
        &\leq \left|\int_{s}^{t_{n+\frac{1}{2}}}(e^{\lambda \theta}L_{1, n+1}u(\theta))''(t_{n+\frac{1}{2}}-\theta)d\theta-\int_{s}^{t_{n}}(e^{\lambda \theta}L_{1, n+1}u(\theta))''(t_{n}-\theta)d\theta\right|\frac{1}{\frac{\tau_{n+1}}{2}}\\
        &=\left|[\int_{s}^{x}(e^{\lambda \theta}L_{1, n+1}u(\theta))''(x-\theta)d\theta]'\right|\\
        &=\left| \int_{s}^{x}(e^{\lambda \theta}L_{1, n+1}u(\theta))''d\theta|_{x=\xi}\right|\\
        &\leq (t_{n+\frac{1}{2}}-t_n)\max_{\theta\in[t_n, t_{n+\frac{1}{2}}]} \left|(e^{\lambda \theta}L_{1, n+1}u(\theta))''\right|
    \end{aligned}
\end{equation}
Then for ${}^{FC}R^{n+\frac{1}{2}}_lu$, we have the following results by (\ref{fe1}), 
\begin{align}
\left| {}^{FC}R^{n+\frac{1}{2}}_lu\right|=&\frac{e^{-\lambda t_{n+\frac{1}{2}}}}{\Gamma(1-\alpha)} \int_{t_{n}}^{t_{n+\frac{1}{2}}} \left\lvert\,  \frac{e^{\lambda t_{n+\frac{1}{2}}} L_{1,  n+1} u\left(t_{n+\frac{1}{2}}\right)-e^{\lambda t_n} L_{1,  n+1} u\left(t_n\right)}{\frac{\tau_{n+1}}{2}}\right.  \notag\\
& \left. -\left(e^{\lambda s} L_{1,  n+1} u(s)\right)^{\prime}\right| \left(t_{n+\frac{1}{2}}-s\right)^{-\alpha} d s\notag\\
\leq& \frac{e^{-\lambda t_{n+\frac{1}{2}}}}{\Gamma(1-\alpha)} \int_{t_{n}}^{t_{n+\frac{1}{2}}} (t_{n+\frac{1}{2}}-t_n)\max_{\theta\in[t_n, t_{n+\frac{1}{2}}]} \left|(e^{\lambda \theta}L_{1, n+1}u(\theta))''\right|\left(t_{n+\frac{1}{2}}-s\right)^{-\alpha}ds\notag\\
=&\frac{e^{-\lambda t_{n+\frac{1}{2}}}}{\Gamma(1-\alpha)} \frac{\tau_{n+1}}{2}\int_{t_{n}}^{t_{n+\frac{1}{2}}} \max_{\theta\in[t_n, t_{n+\frac{1}{2}}]} \left|(e^{\lambda \theta}L_{1, n+1}u(\theta))''\right|\left(t_{n+\frac{1}{2}}-s\right)^{-\alpha}ds\notag\\
=&\frac{e^{-\lambda t_{n+\frac{1}{2}}}}{\Gamma(1-\alpha)} \frac{\tau_{n+1}}{2}\int_{t_{n}}^{t_{n+\frac{1}{2}}} \max_{\theta\in[t_n, t_{n+\frac{1}{2}}]} \left|\lambda^2e^{\lambda \theta}L_{1, n+1}u(\theta)+2\lambda e^{\lambda \theta}\frac{u(t_{n+1})-u(t_n)}{\tau_{n+1}} \right|\left(t_{n+\frac{1}{2}}-s\right)^{-\alpha}ds\notag\\
\leq& \frac{\tau_{n+1}e^{-\lambda t_{n+\frac{1}{2}}}}{2\Gamma(1-\alpha)}\int_{t_{n}}^{t_{n+\frac{1}{2}}} (\lambda^2e^{\lambda t_{n+\frac{1}{2}}}\max_{\theta\in[t_n, t_{n+\frac{1}{2}}]}\left|u(\theta)\right|+2\lambda e^{\lambda t_{n+\frac{1}{2}}}\max_{\theta\in[t_n, t_{n+\frac{1}{2}}]}\left|u'(\theta)\right|)\left(t_{n+\frac{1}{2}}-s\right)^{-\alpha}ds\notag\\
\leq & \frac{\tau_{n+1}}{2\Gamma(1-\alpha)} \int_{t_{n}}^{t_{n+\frac{1}{2}}}(\lambda^2t_{n+\frac{1}{2}}^\delta+2\lambda t_{n+\frac{1}{2}}^{\delta-1})\left(t_{n+\frac{1}{2}}-s\right)^{-\alpha}ds\notag\\
=&\frac{\tau_{n+1}}{2\Gamma(2-\alpha)} (\lambda^2t_{n+\frac{1}{2}}^\delta+2\lambda t_{n+\frac{1}{2}}^{\delta-1})\left(\frac{\tau_{n+\frac{1}{2}}}{2}\right)^{1-\alpha}\notag \\
\leq &K (t_{n+\frac{1}{2}}^{\delta}\tau_{n+1}^{2-\alpha}+t_{n+\frac{1}{2}}^{\delta-1}\tau_{n+1}^{2-\alpha}) \label{0e}\\
\leq &K\left[t_{n+1}^\delta(TN^{-r}n^{r-1})^{2-\alpha}+t_{n+1}^{\delta-1}(TN^{-r}n^{r-1})^{2-\alpha}\right]\notag\\
\leq &K\left[(\frac{n+1}{N})^{r\delta}(\frac{n}{N})^{(r-1)(2-\alpha)}N^{-(2-\alpha)}+(\frac{n+1}{N})^{r(\delta-1)}(\frac{n}{N})^{(r-1)(2-\alpha)}N^{-(2-\alpha)}\right]\notag\\
\leq &KN^{-(2-\alpha)}
\end{align}
Finally, for $n=0, {}^{FC}R^{\frac{1}{2}}_hu=0$ and using (\ref{0e})
    \begin{align}
    |{}^{FC}R^{\frac{1}{2}}_lu|\leq &K (t_{\frac{1}{2}}^{\delta}\tau_{1}^{2-\alpha}+t_{\frac{1}{2}}^{\delta-1}\tau_{1}^{2-\alpha})\notag\\
    \triangleq{}&K(t_1^{2+\delta-\alpha}+t_1^{1+\delta-\alpha})\notag\\
    \triangleq{} &K(N^{-r(2+\delta-\alpha)}+N^{-r(1+\delta-\alpha)})\notag\\
    \leq &KN^{-r(1+\delta-\alpha)}
    \end{align}

We complete the proof by the inequality above. 
\end{proof}


\section{Construction of the full-discrete scheme with fast evaluation}
We first introduce some standard notations. Take the positive integers $M, N$.  Let $h=L/M, t_n=T(\frac{n}{N})^r(0\leq n\leq N)$,  Let's say $\tau_n=t_n-t_{n-1}, x_i=ih(0\leq i\leq M), X_M=\{x_i|0\leq i \leq M\}, T_N=\{t_k|0\leq k\leq N\}$.  Define the mesh function as follows
$$
u_i^n = u(x_i, t_n), 1\leq i \leq M, 1\leq n\leq N. 
$$
$$
f_i^{n+\frac{1}{2}}=f(x_i, t_{n+\frac{1}{2}}), 1\leq i \leq M, 0\leq n\leq N-1. 
$$
For any grid function $u^n=(u_0^n, u_1^n, \cdots, u_M^n)$,  introduce the following notation
\begin{equation}
\begin{aligned}
{\delta}_t u_i^{n+\frac{1}{2}}&=\frac{u_i^{n+1}-u_i^n}{\tau_{n+1}}, \\
\tilde{\delta}_x u_i^n&=\frac{u_{i+1}^n-u_{i-1}^n}{2 h}, \tilde{\delta}_x u_i^{\bar{n}} = \frac{1}{2}(\tilde{\delta}_x u_i^{n+1}+\tilde{\delta}_x u_i^n) \\
\delta_x^2 u_i^n&=\frac{u^n_{i+1}-2u_i^n+u^n_{i-1}}{h^2}, \delta_x^2 u_i^{\bar{n}}=\frac{1}{2}(\delta_x^2 u_i^{n+1}+\delta_x^2 u_i^n). 
\end{aligned}
\end{equation}

Using the fast approximation scheme for Caputo fractional derivatives and using the numerical solution $U_i^n$ instead of the exact solution $u_i^n$,  we get the following Crank-Nicolson difference scheme
\begin{align}
            {\delta_t}U_i^{n+\frac{1}{2}}+{ }_0^{FC}D_t^{\alpha,  \lambda} U_i^{n+\frac{1}{2}}&=\delta_x^2U_i^{\bar{n}}-\tilde{\delta}_x U_i^{\bar{n}}+f_i^{n+\frac{1}{2}}, 1\leq i \leq M-1, 0\leq n \leq N-1, \label{scheme1}\\
        U_i^0&=\phi(x_i), 1\leq i\leq M-1, \label{scheme2}\\
        U_0^n&=U_M^n=0, 0\leq n \leq N\label{scheme3}. 
\end{align}
We can arrange the system above,  so to obtain
\begin{equation}
\left\{\begin{array}{l}
{\left[\left(-\frac{1}{2 h^2}+\frac{1}{4h}\right) U_{i+1}^{n+1}+\left(\frac{1}{\tau_{n+1}}+\frac{1}{2\Gamma(2-\alpha)(\frac{\tau_{n+1}}{2})^\alpha}+\frac{1}{h^2}\right) U_i^{n+1}+\left(-\frac{1}{2 h^2}-\frac{1}{4h}\right) U_{i-1}^{n+1}\right]} \\
+\left[\left(-\frac{1}{2 h^2}+\frac{1}{4h}\right) U_{i+1}^{n}+\left(-\frac{1}{\tau_{n+1}}+\frac{1-2e^{-\lambda \frac{\tau_{n+1}}{2}}}{2\Gamma(2-\alpha)(\frac{\tau_{n+1}}{2})^\alpha}+\frac{1}{h^2}\right) U_i^{n}+\left(-\frac{1}{2 h^2}-\frac{1}{4h}\right) U_{i-1}^{n}\right] \\
-\frac{1}{\Gamma(1-\alpha)}\left[(a_{0, n}-\frac{e^{-\lambda \frac{\tau_{n+1}}{2}}}{(\frac{\tau_{n+1}}{2})^\alpha}) U_i^n+\sum\limits_{l=1}^{n-1} \left(a_{n-l, n}+b_{n-1-l, n}\right) U_i^l+(b_{n-1, n}+\frac{e^{-\lambda t_{n+\frac{1}{2}}}}{t_{n+\frac{1}{2}}^\alpha}) U_i^0\right] \\
=f_i^{n+\frac{1}{2}},  \quad 1 \leq i \leq M-1,  \quad 0 \leq n \leq N-1,  \\
U_i^0=\phi\left(x_i\right),  1 \leq i \leq M-1,  \\
U_0^n=U_M^n=0, 0 \leq n \leq N . 
\end{array}\right. 
\end{equation}
\subsection{The unique solvability of the difference scheme}
\begin{theorem}The difference scheme (\ref{scheme1}-\ref{scheme3}) is uniquely solvable. 
    \begin{proof}
        Let 
        $$
        u^n=(u_0^n, u_1^n, \cdots, u_M^n). 
        $$
        From (\ref{scheme1}-\ref{scheme3}) we know that the value of $U^0$ is known. If the values of the first n layers $U^0, U^1, \cdots, U^{n}$ are uniquely determined. Then a system of linear equations about $U^{n+1}$ can be obtained from (\ref{scheme1}) and (\ref{scheme3}). To prove its unique solvability,  it is only necessary to prove that the corresponding homogeneous system 
        \begin{align}
        &\frac{U_i^{n+1}}{\tau_{n+1}}+\frac{U_i^{n+1}}{2\Gamma(2-\alpha)(\frac{\tau_{n+1}}{2})^\alpha}=\frac{1}{2}\delta_x^2U_i^{n+1}-\frac{1}{2}\tilde{\delta}_xU_i^{n+1}, 1\leq i\leq M-1\notag\\
        &u_0^n=u_M^n=0 \label{unique}
        \end{align}
        has only zero solutions. 

Let $\eta=\frac{1}{\tau_{n+1}}+ \frac{1}{2\Gamma(2-\alpha)(\frac{\tau_{n+1}}{2})^\alpha}$, (\ref{unique}) is equivalent to 
$$
\left[\begin{array}{ccccc}
\eta+\frac{1}{h^2}&\frac{1}{4h}-\frac{1}{2h^2}&0 &\cdots & 0\\
-\frac{1}{4h}-\frac{1}{2h^2} & \eta+\frac{1}{h^2}&\frac{1}{4h}-\frac{1}{2h^2}&\cdots & 0\\
0&-\frac{1}{4h}-\frac{1}{2h^2}&\eta+\frac{1}{h^2}&\cdots&0\\
\vdots & \vdots&\ddots&\ddots & \vdots\\
0&0&\cdots&-\frac{1}{4h}-\frac{1}{2h^2}&\eta+\frac{1}{h^2}
\end{array}
\right]\left(\begin{array}{c}U^{n+1}_1\\
U^{n+1}_2\\
U^{n+1}_3\\
\vdots\\
U^{n+1}_{M-1}
\end{array}\right)\triangleq{}A\left(\begin{array}{c}U^{n+1}_1\\
U^{n+1}_2\\
U^{n+1}_3\\
\vdots\\
U^{n+1}_{M-1}
\end{array}\right)=0
$$
where A is a tridiagonal matrix. And we can get its eigenvalue as follows by \cite{gregory1969}
$$
\lambda_k=\eta+\frac{1}{h^2}+2\sqrt{(\frac{1}{4h}-\frac{1}{2h^2})(-\frac{1}{4h}-\frac{1}{2h^2})}cos(\frac{k\pi}{M}), k=1, 2, \cdots, M-1
$$
Then we know A is invertible because $\lambda_k\geq \eta\geq0$ . So (\ref{unique}) has only zero solutions. 

By using mathematical induction,  we can see that the conclusion of the theorem is established and the theorem has been proved. 
    \end{proof}
\end{theorem}


\subsection{Stability of difference schemes}
\begin{theorem}\label{stable}
    The finite difference scheme (\ref{scheme1})-(\ref{scheme3}) is stable. Let $\{U^n=(U_0^n, U_1^n, \cdots, U_M^n) | 0 \leq n \leq N\}$ and $\{V^n=(V_0^n, V_1^n, \cdots, V_M^n) | 0 \leq n \leq N\}$ be, 
respectively,  the solutions of the problem of (\ref{Equation1})-(\ref{Equation3}) and difference scheme (\ref{scheme1})-(\ref{scheme3}), then 
$$
\left\|U^n-V^n\right\|_{L_2}^2 \leq \left\|U^0-V^0\right\|_{L_2}^2. 
$$
\end{theorem}
Before we give the proof of the theorem, we need some note and lemma first. We analyze the stability of the difference scheme by a Fourier analysis. Define $\rho_j^n=U_j^n-V_j^n, i=0, 1, \cdots, M, n=0, 1, \cdots, N$. Then,  we obtain the following error equation
\begin{equation}\label{errorQ}
\left\{\begin{array}{l}
{\left[\left(-\frac{1}{2 h^2}+\frac{1}{4h}\right) \rho_{j+1}^{n+1}+\left(\frac{1}{\tau_{n+1}}+\frac{1}{2\Gamma(2-\alpha)(\frac{\tau_{n+1}}{2})^\alpha}+\frac{1}{h^2}\right) \rho_j^{n+1}+\left(-\frac{1}{2 h^2}-\frac{1}{4h}\right) \rho_{j-1}^{n+1}\right]} \\
=\left[\left(\frac{1}{2 h^2}-\frac{1}{4h}\right) \rho_{j+1}^{n}+\left(\frac{1}{\tau_{n+1}}-\frac{1-2e^{-\lambda \frac{\tau_{n+1}}{2}}}{2\Gamma(2-\alpha)(\frac{\tau_{n+1}}{2})^\alpha}-\frac{1}{h^2}\right) \rho_j^{n}+\left(\frac{1}{2 h^2}+\frac{1}{4h}\right) \rho_{j-1}^{n}\right] \\
+\frac{1}{\Gamma(1-\alpha)}\left[(a_{0, n}-\frac{e^{-\lambda \frac{\tau_{n+1}}{2}}}{(\frac{\tau_{n+1}}{2})^\alpha}) \rho_j^n+\sum\limits_{l=1}^{n-1} \left(a_{n-l, n}+b_{n-1-l, n}\right) \rho_j^l+(b_{n-1, n}+\frac{e^{-\lambda t_{n+\frac{1}{2}}}}{t_{n+\frac{1}{2}}^\alpha}) \rho_j^0\right] \\
\rho_0^n=\rho_M^n=0, 0 \leq n \leq N . 
\end{array}\right. 
\end{equation}
As the same definition in [13\},  we define the grid function 
\begin{align}
\rho^n(x)= \begin{cases}0,  & 0 \leq x \leq x_{\frac{1}{2}},  \\ \rho_j^n,  & x_{j-\frac{1}{2}} \leq x \leq x_{j+\frac{1}{2}},  \quad 1 \leq j \leq M-1.  \\ 0,  & x_{M-\frac{1}{2}} \leq x \leq x_M . \end{cases}\notag
\end{align}
Since the sequence $\{\rho^n_j\}_{0\leq j\leq M-1}$ is zero at both ends,  it can be extended periodically,  allowing for the application of the Discrete Fourier Transform (DFT) for analysis. 
Denote $\hat{\rho}^n[k]$ as the discrete fourier transform of $\{\rho^n_j\}_{0\leq j\leq M-1}$, then we have
\begin{equation}
\hat{\rho}^n[k]=\sum_{j=0}^{M-1} e^{-i \frac{2 \pi}{M} j k} \rho^n_j \quad k=0, 1,  \ldots,  M-1 . , \notag
\end{equation}
where \( e \) is the base of the natural logarithm,  and \( i \) is the imaginary unit.  The inverse Discrete Fourier Transform (IDFT) is given by:
\[
\rho^n_j = \frac{1}{M} \sum_{k=0}^{M-1} \hat{\rho}^n[k] e^{i \frac{2\pi}{M} jk}
\]
Then we have the Parseval equality for the discrete Fourier transform
$$
\int_0^L\left|\rho^n(x)\right|^2 d x=h\sum_{j=0}^{M-1}\left|\rho^n_j\right|^2=Mh\sum_{k=0}^{M-1}\left|\hat{\rho}^n[k]\right|^2 =L\sum_{k=0}^{M-1}\left|\hat{\rho}^n[k]\right|^2. 
$$
Introduce the following norm
$$
\left\|\rho^n\right\|_{L_2}=\left(\sum_{j=1}^{M-1} h\left|\rho_j^n\right|^2\right)^{1 / 2}=\left(\int_0^L\left|\rho^n\right|^2 d x\right)^{1 / 2}
$$
Then we obtain
$$
\left\|\rho^n\right\|_{L_2}^2=L\sum_{k=0}^{M-1}\left|\hat{\rho}^n[k]\right|^2. 
$$

We multiply both sides of equation (\ref{errorQ}) by $e^{-ijh\beta}$  where $\beta = 2\pi k/M$ and then perform summation, we obtain
\begin{equation}\notag
\begin{aligned}
&\sum_{j=0}^{M-1}\left[\left(-\frac{1}{2 h^2}+\frac{1}{4h}\right) \rho_{j+1}^{n+1}e^{-ijh\beta}+\left(\frac{1}{\tau_{n+1}}+\frac{1}{2\Gamma(2-\alpha)(\frac{\tau_{n+1}}{2})^\alpha}+\frac{1}{h^2}\right) \rho_{j}^{n+1}e^{-ijh\beta}\right. \\
&\left. +\left(-\frac{1}{2 h^2}-\frac{1}{4h}\right) \rho_{j-1}^{n+1}e^{-ijh\beta}\right] \\
&=\sum_{j=0}^{M-1}\left[\left(\frac{1}{2 h^2}-\frac{1}{4h}\right) \rho_{j+1}^{n}e^{-ijh\beta}+\left(\frac{1}{\tau_{n+1}}-\frac{1-2e^{-\lambda \frac{\tau_{n+1}}{2}}}{2\Gamma(2-\alpha)(\frac{\tau_{n+1}}{2})^\alpha}-\frac{1}{h^2}\right) \rho_{j}^{n}e^{-ijh\beta}\right. \\
&\left. +\left(\frac{1}{2 h^2}+\frac{1}{4h}\right) \rho_{j-1}^{n}e^{-ijh\beta}\right] \\
&+\sum_{j=0}^{M-1}\frac{1}{\Gamma(1-\alpha)}\left[(a_{0, n}-\frac{e^{-\lambda \frac{\tau_{n+1}}{2}}}{(\frac{\tau_{n+1}}{2})^\alpha}) \rho_{j}^{n}e^{-ijh\beta}+\sum\limits_{l=1}^{n-1} \left(a_{n-l, n}+b_{n-1-l, n}\right) \rho_{j}^{l}e^{-ijh\beta}\right. \\
&\left. +(b_{n-1, n}+\frac{e^{-\lambda t_{n+\frac{1}{2}}}}{t_{n+\frac{1}{2}}^\alpha}) \rho_{j}^{0}e^{-ijh\beta}\right] 
\end{aligned}
\end{equation}
After simplifications, we get
\begin{equation}\label{simple}
\begin{aligned}
&\left[\left(-\frac{1}{2 h^2}+\frac{1}{4h}\right) e^{ih\beta}+\left(\frac{1}{\tau_{n+1}}+\frac{1}{2\Gamma(2-\alpha)(\frac{\tau_{n+1}}{2})^\alpha}+\frac{1}{h^2}\right)+\left(-\frac{1}{2 h^2}-\frac{1}{4h}\right) e^{-ih\beta}\right]\hat{\rho}^{n+1}[k] \\
&=\left[\left(\frac{1}{2 h^2}-\frac{1}{4h}\right) e^{ih\beta}+\left(\frac{1}{\tau_{n+1}}-\frac{1-2e^{-\lambda \frac{\tau_{n+1}}{2}}}{2\Gamma(2-\alpha)(\frac{\tau_{n+1}}{2})^\alpha}-\frac{1}{h^2}\right)+\left(\frac{1}{2 h^2}+\frac{1}{4h}\right) e^{-ih\beta}\right]\hat{\rho}^{n}[k] \\
&+\frac{1}{\Gamma(1-\alpha)}\left[(a_{0, n}-\frac{e^{-\lambda \frac{\tau_{n+1}}{2}}}{(\frac{\tau_{n+1}}{2})^\alpha}) \hat{\rho}^{n}[k]+\sum\limits_{l=1}^{n-1} \left(a_{n-l, n}+b_{n-1-l, n}\right) \hat{\rho}^{l}[k]+(b_{n-1, n}+\frac{e^{-\lambda t_{n+\frac{1}{2}}}}{t_{n+\frac{1}{2}}^\alpha}) \hat{\rho}^{0}[k]\right] 
\end{aligned}
\end{equation}
Noting that $e^{ih\beta}+e^{-ih\beta}=2cos(h\beta)$ and $e^{ih\beta}-e^{-ih\beta}=2isin(h\beta)$, thus we can simplify (\ref{simple}) into the following fomulation
\begin{align}
    &\left[-\frac{1}{h^2}cos(h\beta)+\frac{1}{2h}isin(h\beta)+\left(\frac{1}{\tau_{n+1}}+\frac{1}{2\Gamma(2-\alpha)(\frac{\tau_{n+1}}{2})^\alpha}+\frac{1}{h^2}\right) \right]\hat{\rho}^{n+1}[k]\notag \\
    &=\left[\frac{1}{h^2}cos(h\beta)-\frac{1}{2h}isin(h\beta)+\left(\frac{1}{\tau_{n+1}}-\frac{1-2e^{-\lambda \frac{\tau_{n+1}}{2}}}{2\Gamma(2-\alpha)(\frac{\tau_{n+1}}{2})^\alpha}-\frac{1}{h^2}\right) -\frac{e^{-\lambda \frac{\tau_{n+1}}{2}}}{\Gamma(1-\alpha)(\frac{\tau_{n+1}}{2})^\alpha}\right]\hat{\rho}^{n}[k] \notag\\
    &+\frac{1}{\Gamma(1-\alpha)}\left[a_{0, n} \hat{\rho}^{n}[k]+\sum\limits_{l=1}^{n-1} \left(a_{n-l, n}+b_{n-1-l, n}\right) \hat{\rho}^{l}[k]+(b_{n-1, n}+\frac{e^{-\lambda t_{n+\frac{1}{2}}}}{t_{n+\frac{1}{2}}^\alpha}) \hat{\rho}^{0}[k]\right]\notag\\
    &=\left[\frac{1}{h^2}cos(h\beta)-\frac{1}{2h}isin(h\beta)+\left(\frac{1}{\tau_{n+1}}-\frac{1}{2\Gamma(2-\alpha)(\frac{\tau_{n+1}}{2})^\alpha}-\frac{1}{h^2}\right) \right]\hat{\rho}^{n}[k] \notag\\
    &+\frac{1}{\Gamma(1-\alpha)}\left[(a_{0, n} +\frac{e^{-\lambda \frac{\tau_{n+1}}{2}}}{(1-\alpha)(\frac{\tau_{n+1}}{2})^\alpha}-\frac{e^{-\lambda \frac{\tau_{n+1}}{2}}}{(\frac{\tau_{n+1}}{2})^\alpha})\hat{\rho}^{n}[k]\right.\notag\\
    &+\left.\sum\limits_{l=1}^{n-1} \left(a_{n-l, n}+b_{n-1-l, n}\right) \hat{\rho}^{l}[k]+(b_{n-1, n}+\frac{e^{-\lambda t_{n+\frac{1}{2}}}}{t_{n+\frac{1}{2}}^\alpha}) \hat{\rho}^{0}[k]\right]\label{d0}
\end{align}

\begin{lemma}\label{dnd0}
    Suppose that $\hat{\rho}_{n}[k] (n = 1,  2, \cdots,  N)$ are solutions of (\ref{d0}).  Then we have
    $$
    |\hat{\rho}^{n}[k]|\leq |\hat{\rho}^{0}[k]|, n = 1,  2, \cdots,  N
    $$
\end{lemma}
\begin{proof}
When $n=0$, we get
\begin{align}
    |\hat{\rho}_{1}[k]|&=\left|\frac{\frac{1}{h^2}cos(h\beta)-\frac{1}{2h}isin(h\beta)+\left(\frac{1}{\tau_{1}}-\frac{1-2e^{-\lambda \frac{\tau_{1}}{2}}}{2\Gamma(2-\alpha)(\frac{\tau_{1}}{2})^\alpha}-\frac{1}{h^2}\right) }{-\frac{1}{h^2}cos(h\beta)+\frac{1}{2h}isin(h\beta)+\left(\frac{1}{\tau_{1}}+\frac{1}{2\Gamma(2-\alpha)(\frac{\tau_{1}}{2})^\alpha}+\frac{1}{h^2}\right) }\right||\hat{\rho}^{0}[k]| \notag\\
    &=\sqrt{\frac{\left( \frac{1}{\tau_1}-\frac{1-2e^{-\lambda \frac{\tau_{1}}{2}}}{2\Gamma(2-\alpha)(\frac{\tau_{1}}{2})^\alpha}-\frac{1-cos(h\beta)}{h^2}\right)^2+\left(\frac{sin(h\beta)}{2h}\right)^2}{\left( \frac{1}{\tau_1}+\frac{1}{2\Gamma(2-\alpha)(\frac{\tau_{1}}{2})^\alpha}+\frac{1-cos(h\beta)}{h^2}\right)^2+\left(\frac{sin(h\beta)}{2h}\right)^2}}|\hat{\rho}^{0}[k]| \notag\\
    &\leq |\hat{\rho}^{0}[k]|
\end{align}
Then we suppose $|\hat{\rho}^{m}[k]|\leq |\hat{\rho}^{0}[k]|, m=1, 2, \cdots, n$. 
\begin{align}
    &\left|\frac{1}{\Gamma(1-\alpha)}\left[(a_{0, n} +\frac{e^{-\lambda \frac{\tau_{n+1}}{2}}}{(1-\alpha)(\frac{\tau_{n+1}}{2})^\alpha}-\frac{e^{-\lambda \frac{\tau_{n+1}}{2}}}{(\frac{\tau_{n+1}}{2})^\alpha})\hat{\rho}^{n}[k]\right.\right.\notag\\
    &+\left.\left.\sum\limits_{l=1}^{n-1} \left(a_{n-l, n}+b_{n-1-l, n}\right) \hat{\rho}^{l}[k]+(b_{n-1, n}+\frac{e^{-\lambda t_{n+\frac{1}{2}}}}{t_{n+\frac{1}{2}}^\alpha}) \hat{\rho}^{0}[k]\right]\right|\notag\\
    &\leq \left[\frac{e^{-\lambda \frac{\tau_{n+1}}{2}}}{\Gamma(2-\alpha)(\frac{\tau_{n+1}}{2})^\alpha}-\frac{e^{-\lambda \frac{\tau_{n+1}}{2}}}{\Gamma(1-\alpha)(\frac{\tau_{n+1}}{2})^\alpha}\right]|\hat{\rho}^{0}[k]|+ \notag\\
    &\frac{1}{\Gamma(1-\alpha)}\sum_{l=0}^{n-1}(a_{l, n}+b_{l, n})|\hat{\rho}^{0}[k]|+\frac{1}{\Gamma(1-\alpha)}\frac{e^{-\lambda t_{n+\frac{1}{2}}}}{t_{n+\frac{1}{2}}^\alpha} |\hat{\rho}^{0}[k]| \label{dn}
\end{align}
We have the following results by Lemma \ref{fast}, (\ref{lambda}) and (\ref{ab})
\begin{align}
  \sum_{l=0}^{n-1}(a_{l, n}+b_{l, n})  &=\sum_{l=0}^{n-1}\alpha \sum_{i=1}^{N_{exp}}w_ie^{-(\lambda+s_i)(t_{n+\frac{1}{2}}-t_{n+\frac{1}{2}-l})}\int_{t_{n-l-1}}^{t_{n-l}}e^{-(\lambda+s_i)(t_{n+\frac{1}{2}-l}-s)}ds\notag\\
  &=\sum_{l=0}^{n-1}\alpha \sum_{i=1}^{N_{exp}}w_i\int_{t_{n-l-1}}^{t_{n-l}}e^{-(\lambda+s_i)(t_{n+\frac{1}{2}}-s)}ds\notag\\
  &=\sum_{l=0}^{n-1}\alpha\int_{t_{n-l-1}}^{t_{n-l}}\sum_{i=1}^{N_{exp}}w_ie^{-s_i(t_{n+\frac{1}{2}}-s)}e^{-\lambda(t_{n+\frac{1}{2}}-s)}ds\notag\\
  &=\sum_{l=0}^{n-1}\alpha\int_{t_{n-l-1}}^{t_{n-l}}\frac{1}{(t_{n+\frac{1}{2}}-s)^{1+\alpha}}e^{-\lambda(t_{n+\frac{1}{2}}-s)}ds\notag\\
  &=\alpha\int_{t_{0}}^{t_{n}}\frac{1}{(t_{n+\frac{1}{2}}-s)^{1+\alpha}}e^{-\lambda(t_{n+\frac{1}{2}}-s)}ds\notag\\
  &=\alpha\int_{\frac{\tau_{n+1}}{2}}^{t_{n+\frac{1}{2}}}\frac{e^{-\lambda s}}{s^{1+\alpha}}ds\notag\\
  &=\frac{e^{-\lambda \frac{\tau_{n+1}}{2}}}{(\frac{\tau_{n+1}}{2})^\alpha}-\frac{e^{-\lambda t_{n+\frac{1}{2}}}}{t_{n+\frac{1}{2}}^\alpha}-\lambda\int_{\frac{\tau_{n+1}}{2}}^{t_{n+\frac{1}{2}}}\frac{e^{-\lambda s}}{s^{\alpha}}ds\notag\\
  &\leq \frac{e^{-\lambda \frac{\tau_{n+1}}{2}}}{(\frac{\tau_{n+1}}{2})^\alpha}-\frac{e^{-\lambda t_{n+\frac{1}{2}}}}{t_{n+\frac{1}{2}}^\alpha}\label{aplusb}
\end{align}
Combine \ref{d0},\ref{dn} and \ref{aplusb}, we have
\begin{align}
    &\left|-\frac{1}{h^2}cos(h\beta)+\frac{1}{2h}isin(h\beta)+\left(\frac{1}{\tau_{n+1}}+\frac{1}{2\Gamma(2-\alpha)(\frac{\tau_{n+1}}{2})^\alpha}+\frac{1}{h^2}\right) \right|\left|\hat{\rho}^{n+1}[k]\right|\notag \\
    \leq&\left|\frac{1}{h^2}cos(h\beta)-\frac{1}{2h}isin(h\beta)+\left(\frac{1}{\tau_{n+1}}-\frac{1}{2\Gamma(2-\alpha)(\frac{\tau_{n+1}}{2})^\alpha}-\frac{1}{h^2}\right) \right|\left|\hat{\rho}^{n}[k]\right| \notag\\
    &+\frac{e^{-\lambda \frac{\tau_{n+1}}{2}}}{\Gamma(2-\alpha)(\frac{\tau_{n+1}}{2})^\alpha}\left|\hat{\rho}^{0}[k]\right|\label{need1}
\end{align}
Let
\begin{align*}
A &= \frac{1-cos(h\beta)}{h^2},B=\frac{1}{2\Gamma(2-\alpha)(\frac{\tau_{n+1}}{2})^\alpha},C=2Be^{-\lambda\frac{\tau_{n+1}}{2}};\\
a&=-A+\frac{1}{\tau_{n+1}}-B,b=A+\frac{1}{\tau_{n+1}}+B,d=\frac{sin(h\beta)}{2h},
\end{align*}
then we want to proof the following inequality
\begin{align*}
&\sqrt{a^2+d^2}+C\leq\sqrt{b^2+d^2}\\
\iff&\sqrt{a^2+d^2}+\sqrt{b^2+d^2}\leq \frac{(b-a)(b+a)}{C}\\
\iff&a^2+d^2\leq \frac{(b-a)^2(b+a)^2}{C^2}-\frac{2(b-a)(b+a)}{C}\sqrt{b^2+d^2}+b^2+d^2\\
\iff&\sqrt{b^2+d^2}\leq \frac{(b-a)(b+a)}{2C}+\frac{C}{2}\\
\iff&(A+B)^2+\frac{1}{\tau_{n+1}^2}+d^2\leq B^2e^{-\lambda\tau_{n+1}}+\frac{(A+B)^2}{B^2e^{-\lambda\tau_{n+1}}\tau_{n+1}^2}\\
\iff&d^2\leq \frac{1-B^2e^{-\lambda\tau_{n+1}}\tau_{n+1}^2}{B^2e^{-\lambda\tau_{n+1}}\tau_{n+1}^2}((A+B)^2-B^2e^{-\lambda\tau_{n+1}})
\end{align*}
noticing that $d^2 = \frac{sin(h\beta)^2}{4h^2}\leq \frac{1-cos(h\beta)^2}{4h^2}\leq \frac{A}{2}$ and $(A+B)^2-B^2e^{-\lambda\tau_{n+1}}\geq A^2+2AB$, so $\frac{(A+B)^2-B^2e^{-\lambda\tau_{n+1}}}{d^2}\geq 2\frac{A^2+2AB}{A}\geq 4B$ which means we only need
\begin{align}
    1\leq 4B\frac{1-B^2e^{-\lambda\tau_{n+1}}\tau_{n+1}^2}{B^2e^{-\lambda\tau_{n+1}}\tau_{n+1}^2}\label{need2}
\end{align}
Combining the obtained expression with the condition $\frac{1}{2}<\Gamma(2-\alpha)<1$, we employ bounding arguments to simplify the inequality \ref{need2}, and it is therefore sufficient to focus on the relation $(\frac{\tau_{n+1}}{2})^{2-2\alpha}<\frac{1}{3}$ which is roughly estimated and easily ensured in numerical simulation. Finally, we get $|\hat{\rho}^{n+1}[k]|\leq |\hat{\rho}^{0}[k]|$ by equation (\ref{need1}) and (\ref{need2}).  By using mathematical induction,  we complete the proof. 
\end{proof}
Now we give the proof of theorem \ref{stable}. 
\begin{proof}
    Using Lemma \ref{dnd0} and the Parseval equality,  we obtain
    \begin{equation}
\begin{aligned}
\left\|U^n-V^n\right\|_{L_2}^2 & =\left\|\rho^n\right\|_{L_2}^2=\sum_{j=1}^{M-1} h\left|\rho_j^n\right|^2=L\sum_{k=0}^{M-1}|\hat{\rho}^n[k]|^2 \\
& \leq L\sum_{k=0}^{M-1}|\hat{\rho}^0[k]|^2=\left\|\rho^0\right\|_{L_2}^2=\left\|U^0-V^0\right\|_{L_2}^2, 
\end{aligned}
\end{equation}
which prove the scheme is stable. 
\end{proof}

\subsection{Convergence of difference schemes}
Now,  we consider the convergence of the
difference scheme. 
\begin{theorem}\label{converge}
    Suppose that $u^n$ is the exact solution of problem (\ref{Equation1})–(\ref{Equation3}) with $\left|\frac{\partial^l u(x, t)}{\partial t^l}\right| \leq C(1+t^{\delta-l}),  l=0, 1, 2, 3, \left|\frac{\partial^k u(x, t)}{\partial x^k}\right| \leq C, k=0, 1, 2, 3, 4, \left|\frac{\partial^{2+p} u(x, t)}{\partial t^2\partial x^p}\right| \leq C(1+t^{\delta-2}), p=1, 2$ and $r(\delta-1) > 2, r\geq 3$.  Then the numerical solution  $U^n(n = 0,  1,  \cdots ,  N)$ to (\ref{scheme1})-(\ref{scheme3}) satisfies
    $$
    ||u^n-U^n||_{L_2}^2\leq C(N^{-2}+h^2+\varepsilon)^2
    $$
\end{theorem}
Although there are two conditions involving $r$ in Theorem \ref{converge},  since $ r $ can be adjusted arbitrarily,  we can set \( r \geq \max\left\{\frac{2}{\delta - 1},  3\right\} \) once \( \delta \) and \( \alpha \) are given.  In order to prove the convergence theorem,  we first introduce several lemmas. 
\begin{lemma} \label{second}
Suppose that $u(x, t)$ satisfies $\left|\frac{\partial^l u(x, t)}{\partial t^l}\right| \leq C(1+t^{\delta-l}),  l=2, 3, \left|\frac{\partial^k u(x, t)}{\partial x^k}\right| \leq C, k=0, 1, 2, 3, 4, \left|\frac{\partial^{2+p} u(x, t)}{\partial t^2\partial x^p}\right| \leq C(1+t^{\delta-2}), p=1, 2$ , then there exists a constant K such that
\begin{equation}
    \begin{aligned}
    |u_t\left(x_i,  t_{n+\frac{1}{2}}\right) -{\delta}_t u_i^{n+\frac{1}{2}}| &\triangleq{}|R_t^{i, n+\frac{1}{2}}|\leq K(n+1)^{-2}\\
    |u_x(x_i, t_{n+\frac{1}{2}})-\tilde{\delta}_x u_i^{\bar{n}}|&\triangleq{}R_x^{i, n+\frac{1}{2}}\leq K[(n+1)^{-2}+h^{2}]\\
    |u_{xx}(x_i, t_{n+\frac{1}{2}})-\delta_x^2u_i^{\bar{n}}|&\triangleq{}R_{xx}^{i, n+\frac{1}{2}}\leq K[(n+1)^{-2}+h^{2}]
    \end{aligned}
\end{equation}
especially for $n=0$, we have a higher progress estimate
\begin{equation}
\begin{aligned}
|u_t\left(x_i,  t_{\frac{1}{2}}\right) -{\delta}_t u_i^{n+\frac{1}{2}}| &\triangleq{}|R_t^{i, n+\frac{1}{2}}|\leq KN^{-r(\delta-1)}\\
    |u_x(x_i, t_{\frac{1}{2}})-\tilde{\delta}_x u_i^{\bar{n}}|&\triangleq{}R_x^{i, n+\frac{1}{2}}\leq K[N^{-r\delta}+h^{2}]\\
    |u_{xx}(x_i, t_{\frac{1}{2}})-\delta_x^2u_i^{\bar{n}}|&\triangleq{}R_{xx}^{i, n+\frac{1}{2}}\leq K[N^{-r\delta}+h^{2}] \label{spaceerror}\end{aligned}
\end{equation}
\end{lemma}
\begin{proof}
    Firstly, for $n\geq 0$, we can use the Taylor expansion and get
    \begin{align}
        u_i^{n+1}&=u_i^{n+\frac{1}{2}}+u_t(x_i, t_{n+\frac{1}{2}})(t_{n+1}-t_{n+\frac{1}{2}})+\frac{1}{2}u_{tt}(x_i, t_{n+\frac{1}{2}})(t_{n+1}-t_{n+\frac{1}{2}})^2\notag\\
        &+\frac{1}{2}\int_{t_{n+\frac{1}{2}}}^{t_{n+1}}\frac{\partial^3 u}{\partial t^3}(t_{n+1}-t)^2dt \notag\\
        u_i^{n}&=u_i^{n+\frac{1}{2}}+u_t(x_i, t_{n+\frac{1}{2}})(t_{n}-t_{n+\frac{1}{2}})++\frac{1}{2}u_{tt}(x_i, t_{n+\frac{1}{2}})(t_{n}-t_{n+\frac{1}{2}})^2\notag\\
        &+\frac{1}{2}\int_{t_{n+\frac{1}{2}}}^{t_{n}}\frac{\partial^3 u}{\partial t^3}(t_{n}-t)^2dt \notag
    \end{align}
    then for $n\geq 1$ we get 
    \begin{align}
\left|u_t(x_i, t_{n+\frac{1}{2}})-\delta_tu_i^{n+\frac{1}{2}}\right| & =\left|u_t(x_i, t_{n+\frac{1}{2}})-\frac{u_i^{n+1}-u_i^n}{\tau_{n+1}}\right| \notag\\
& =\left|\frac{\frac{1}{2}\int_{t_{n+\frac{1}{2}}}^{t_{n+1}}\frac{\partial^3 u}{\partial t^3}(t_{n+1}-t)^2dt-\frac{1}{2}\int_{t_{n+\frac{1}{2}}}^{t_{n}}\frac{\partial^3 u}{\partial t^3}(t_{n}-t)^2dt}{\tau_{n+1}}\right| \notag\\
& =\left|\left[\frac{1}{2}\int_{t_{n+\frac{1}{2}}}^x \frac{\partial^3 u}{\partial t^3}(x-t)^2 d t\right]^{\prime}\right|_{x=\xi} \label{mean} \\
& =\left|\left[\int_{t_{n+\frac{1}{2}}}^x \frac{\partial^3 u}{\partial t^3}(x-t) d t\right]\right|_{x=\xi}  \notag\\
& =\left|\int_{t_{n+\frac{1}{2}}}^{\xi} \frac{\partial^3 u}{\partial t^3}(\xi - t) d t\right|,  \xi \in\left[t_n,  t_{n+1}\right] \notag\\
& \leq \frac{\tau_{n+1}^2}{2} \max _{t \in\left[t_n,  t_{n+1}\right]}\left|\frac{\partial^3 u}{\partial t^3}\right| \notag\\
& \leq C \tau_{n+1}^2 t_n^{\delta-3} \notag\\
&\leq CTN^{-2r}n^{2(r-1)}[T(\frac{n}{N})^r]^{\delta-3} \label{uselemma}\\
&\triangleq{}K\frac{n^{r\delta-r-2}}{N^{r\delta-r}}\notag\\
&\leq Kn^{-2}\notag
\end{align}
where (\ref{uselemma}) is obtained by Lemma \ref{tau} . And for $n=0$ we have
\begin{align}
\left|u_t(x_i, t_{\frac{1}{2}})-\delta_tu_i^{\frac{1}{2}}\right| & =\left|u_t(x_i, t_{\frac{1}{2}})-\frac{u_i^{1}-u_i^0}{\tau_{1}}\right| \notag\\
& =\left|\frac{\frac{1}{2}\int_{t_{\frac{1}{2}}}^{t_{1}}\frac{\partial^3 u}{\partial t^3}(t_{1}-t)^2dt-\frac{1}{2}\int_{t_{\frac{1}{2}}}^{t_{0}}\frac{\partial^3 u}{\partial t^3}(t_{0}-t)^2dt}{\tau_{1}}\right| \notag\\
&=\left|\frac{\int_{t_{\frac{1}{2}}}^{t_{1}}\frac{\partial^3 u}{\partial t^3}(t_{1}^2-2t_1t)dt+\int_{t_{\frac{1}{2}}}^{t_{1}}\frac{\partial^3 u}{\partial t^3}t^2dt-\int_{t_{\frac{1}{2}}}^{t_{0}}\frac{\partial^3 u}{\partial t^3}t^2dt}{2\tau_{1}}\right| \notag\\
&=\left|\frac{\int_{t_{\frac{1}{2}}}^{t_{1}}\frac{\partial^3 u}{\partial t^3}(t_{1}^2-2t_1t)dt+\int_{t_{0}}^{t_{1}}\frac{\partial^3 u}{\partial t^3}t^2dt}{2\tau_{1}}\right| \notag\\
&\leq \frac{\left|\int_{t_{\frac{1}{2}}}^{t_{1}}\frac{\partial^3 u}{\partial t^3}(t_{1}^2-2t_1t)dt\right|+\left|\int_{t_{0}}^{t_{1}}\frac{\partial^3 u}{\partial t^3}t^2dt\right|}{2\tau_{1}} \notag\\
&\leq \frac{t_1^2\int_{t_{\frac{1}{2}}}^{t_{1}}|\frac{\partial^3 u}{\partial t^3}|dt+2t_1\int_{t_{\frac{1}{2}}}^{t_{1}}|\frac{\partial^3 u}{\partial t^3}|tdt+\int_{t_{0}}^{t_{1}}|\frac{\partial^3 u}{\partial t^3}|t^2dt}{2\tau_1}\notag\\
&\leq \frac{t_1^2\int_{t_{\frac{1}{2}}}^{t_{1}}Ct^{\delta-3}dt+2t_1\int_{t_{\frac{1}{2}}}^{t_{1}}Ct^{\delta-2}dt+\int_{t_{0}}^{t_{1}}Ct^{\delta-1}dt}{2t_1}\notag\\
&=[t_1t_{\frac{1}{2}}^{\delta-2}+2t_{\frac{1}{2}}^{\delta-1}+\frac{t_1^{\delta-1}}{\delta}]\notag\\
&= C(\frac{T}{N^r})^{\delta-1}\notag\\
&\leq KN^{-r(\delta-1)}\notag
\end{align}
so $|R_t^{i, n+\frac{1}{2}}|\leq K(n+1)^{-2}$. 

Secondly, we will estimate $R_x^{i, n+\frac{1}{2}}$, also by the Taylor expansion we have
$$
u_x(x_i, t_n)=\frac{u_{i+1}^n-u_{i-1}^n}{2h}+O(h^2)=\tilde{\delta}_x u_i^n+O(h^2), 
$$
$$
u_x(x_i, t_{n+1})=\frac{u_{i+1}^{n+1}-u_{i-1}^{n+1}}{2h}+O(h^2)=\tilde{\delta}_x u_i^{n+1}+O(h^2), 
$$
$$
u_x(x_i, t_{n+1})=u_x(x_i, t_{n+\frac{1}{2}})+u_{xt}(x_i, t_{n+\frac{1}{2}})\frac{\tau_{n+1}}{2}+\int_{t_{n+\frac{1}{2}}}^{t_{n+1}}u_{xtt}(x_i, t)(t_{n+1}-t)dt, 
$$
$$
u_x(x_i, t_{n})=u_x(x_i, t_{n+\frac{1}{2}})-u_{xt}(x_i, t_{n+\frac{1}{2}})\frac{\tau_{n+1}}{2}+\int_{t_{n+\frac{1}{2}}}^{t_{n}}u_{xtt}(x_i, t)(t_{n}-t)dt
$$
From the equations above and use integration by parts, we have
\begin{align}
    |u_x(x_i, t_{n+\frac{1}{2}})-\tilde{\delta}_xu_i^{\bar{n}}|&\leq |\frac{\int_{t_{n+\frac{1}{2}}}^{t_{n+1}}u_{xtt}(x_i, t)(t_{n+1}-t)dt+\int_{t_{n+\frac{1}{2}}}^{t_{n}}u_{xtt}(x_i, t)(t_{n}-t)dt}{2}|+O(h^2)\notag\\
    &\leq \frac{Ct_{n+\frac{1}{2}}^{\delta-2}\int_{t_{n+\frac{1}{2}}}^{t_{n+1}}(t_{n+1}-t)dt+C\int^{t_{n+\frac{1}{2}}}_{t_{n}}t^{\delta-2}(t_{n}-t)dt}{2}+O(h^2)\notag\\
    &\leq C(t_{n+1}^{\delta-2}\tau_{n+1}^2+\int^{t_{n+\frac{1}{2}}}_{t_{n}}t^{\delta-2}(t-t_{n})dt)+O(h^2)\label{yijiechashang}
\end{align}
Similarly, we consider $n=0$ and $n\geq 1$, when $n=0$
\begin{align*}
    |u_x(x_i, t_{n+\frac{1}{2}})-\tilde{\delta}_xu_i^{\bar{n}}|&\leq C(t_1^\delta+\int_0^{t_{\frac{1}{2}}}t^{\delta-1}dt)+O(h^2)\\
    &\leq 2Ct_1^\delta+O(h^2)\\
    &=2CT^\delta\frac{1}{N^{r\delta}}+O(h^2)\\
    &\leq KN^{-r\delta}+O(h^2)
\end{align*}
when $n\geq 1$, using $t_{n+1}\leq 2^rt_n$ and we have
\begin{align*}
    |u_x(x_i, t_{n+\frac{1}{2}})-\tilde{\delta}_xu_i^{\bar{n}}|&\leq C(t_{n+1}^{\delta-2}\tau_{n+1}^2+t_n^{\delta-2}\int^{t_{n+\frac{1}{2}}}_{t_{n}}(t-t_{n})dt)+O(h^2)\\
    &\leq C(t_{n+1}^{\delta-2}\tau_{n+1}^2+t_n^{\delta-2}\frac{\tau_{n+1}^2}{8})+O(h^2)\\
    &\leq Kt_{n+1}^{\delta-2}\tau_{n+1}^2+O(h^2)\\
    &\triangleq{} K(\frac{n+1}{N})^{r(\delta-2)}\frac{(n+1)^{2(r-1)}}{N^{2r}}+O(h^2)\\
    &\leq K\frac{(n+1)^{r\delta-2}}{N^{r\delta}}+O(h^2)\\
    &\leq K(n+1)^{-2}+O(h^2)
\end{align*}
so $|R_x^{i, n+\frac{1}{2}}|\leq K[(n+1)^{-2}+h^2]$ . 

Finally,  the proof of $R_{xx}^{i, n+\frac{1}{2}}$ is similar with $R_x^{i, n+\frac{1}{2}}$, also by the Taylor expansion we have
$$
u_{xx}(x_i, t_n)=\frac{u_{i+1}^n-2u_i^n+u_{i-1}^n}{h^2}+O(h^2)=\delta_x^2u_i^n+O(h^2)
$$
$$
u_{xx}(x_i, t_{n+1})=\frac{u_{i+1}^{n+1}-2u_i^{n+1}+u_{i-1}^{n+1}}{h^2}+O(h^2)=\delta_x^2u_i^{n+1}+O(h^2)
$$
$$
u_{xx}(x_i, t_{n+1})=u_{xx}(x_i, t_{n+\frac{1}{2}})+u_{xxt}(x_i, t_{n+\frac{1}{2}})\frac{\tau_{n+1}}{2}+\int_{t_{n+\frac{1}{2}}}^{t_{n+1}}u_{xxtt}(x_i, t)(t_{n+1}-t)dt, 
$$
$$
u_{xx}(x_i, t_{n})=u_{xx}(x_i, t_{n+\frac{1}{2}})-u_{xxt}(x_i, t_{n+\frac{1}{2}})\frac{\tau_{n+1}}{2}+\int_{t_{n+\frac{1}{2}}}^{t_{n}}u_{xxtt}(x_i, t)(t_{n}-t)dt
$$
From the equations above and use integration by parts, we have
\begin{align}
    |u_{xx}(x_i, t_{n+\frac{1}{2}})-\delta_x^2u_i^{\bar{n}}|&\leq |\frac{\int_{t_{n+\frac{1}{2}}}^{t_{n+1}}u_{xxtt}(x_i, t)(t_{n+1}-t)dt+\int_{t_{n+\frac{1}{2}}}^{t_{n}}u_{xxtt}(x_i, t)(t_{n}-t)dt}{2}|+O(h^2)\notag\\
    &\leq \frac{Ct_{n+\frac{1}{2}}^{\delta-2}\int_{t_{n+\frac{1}{2}}}^{t_{n+1}}(t_{n+1}-t)dt+C\int^{t_{n+\frac{1}{2}}}_{t_{n}}t^{\delta-2}(t_{n}-t)dt}{2}+O(h^2)\notag\\
    &\leq C(t_{n+1}^{\delta-2}\tau_{n+1}^2+\int^{t_{n+\frac{1}{2}}}_{t_{n}}t^{\delta-2}(t-t_{n})dt)+O(h^2)\notag
\end{align}
which is exactly same with the inequation (\ref{yijiechashang}), so we finish the proof. 
\end{proof}
With the above lemma,  we can proceed with the following analysis. Let $u_i^n$ be the solution of problem (\ref{Equation1})-(\ref{Equation3}) at $(x_i, t_n)$, From Lemma \ref{FRu} and Lemma \ref{second}, we have
\begin{equation}
\left\{\begin{array}{l}
{\left[\left(-\frac{1}{2 h^2}+\frac{1}{4h}\right) u_{i+1}^{n+1}+\left(\frac{1}{\tau_{n+1}}+\frac{1}{2\Gamma(2-\alpha)(\frac{\tau_{n+1}}{2})^\alpha}+\frac{1}{h^2}\right) u_i^{n+1}+\left(-\frac{1}{2 h^2}-\frac{1}{4h}\right) u_{i-1}^{n+1}\right]} \\
+\left[\left(-\frac{1}{2 h^2}+\frac{1}{4h}\right) u_{i+1}^{n}+\left(-\frac{1}{\tau_{n+1}}+\frac{1-2e^{-\lambda \frac{\tau_{n+1}}{2}}}{2\Gamma(2-\alpha)(\frac{\tau_{n+1}}{2})^\alpha}+\frac{1}{h^2}\right) u_i^{n}+\left(-\frac{1}{2 h^2}-\frac{1}{4h}\right) u_{i-1}^{n}\right] \\
-\frac{1}{\Gamma(1-\alpha)}\left[(a_{0, n}-\frac{e^{-\lambda \frac{\tau_{n+1}}{2}}}{(\frac{\tau_{n+1}}{2})^\alpha}) u_i^n+\sum\limits_{l=1}^{n-1} \left(a_{n-l, n}+b_{n-1-l, n}\right) u_i^l+(b_{n-1, n}+\frac{e^{-\lambda t_{n+\frac{1}{2}}}}{t_{n+\frac{1}{2}}^\alpha}) u_i^0\right] \\
=f_i^{n+\frac{1}{2}}+{}^{FC}R^{n+\frac{1}{2}}+R_t^{i, n+\frac{1}{2}}+R_x^{i, n+\frac{1}{2}}+R_{xx}^{i, n+\frac{1}{2}},  \quad 1 \leq i \leq M-1,  \quad 0 \leq n \leq N-1,  \\
u_i^0=\phi\left(x_i\right),  1 \leq i \leq M-1,  \\
u_0^n=u_M^n=0, 0 \leq n \leq N . 
\end{array}\right. 
\end{equation}
Denote the truncation error at $(t_{n+\frac{1}{2}}, x_j)$ as $R_j^{n+\frac{1}{2}}$. Then we know $R_j^{n+\frac{1}{2}}={}^{FC}R^{n+\frac{1}{2}}+R_t^{i, n+\frac{1}{2}}+R_x^{i, n+\frac{1}{2}}+R_{xx}^{i, n+\frac{1}{2}}$ and 
$$
|R_j^{n+\frac{1}{2}}|\leq K[(n+1)^{-(2-\alpha)}+h^2+\varepsilon], n\geq 0
$$
Particularly for n=0, we have the following estimation by equation (\ref{L1half}) ,  (\ref{Fhalf}) , (\ref{spaceerror})and $r\geq \frac{2}{\delta-1}$
\begin{align}
|R_j^{\frac{1}{2}}|\leq K[N^{-min[r(1+\delta-\alpha), -r(\delta-1), -r\delta]}+h^2+\varepsilon]\leq K[N^{-2}+h^2+\varepsilon] \label{R_half}
\end{align}
Let $\varepsilon_i^n=u_i^n-U_i^n$ and we can get the error equation 
\begin{align}
&{\left[\left(-\frac{1}{2 h^2}+\frac{1}{4h}\right) \varepsilon_{j+1}^{n+1}+\left(\frac{1}{\tau_{n+1}}+\frac{1}{2\Gamma(2-\alpha)(\frac{\tau_{n+1}}{2})^\alpha}+\frac{1}{h^2}\right) \varepsilon_j^{n+1}+\left(-\frac{1}{2 h^2}-\frac{1}{4h}\right) \varepsilon_{j-1}^{n+1}\right]} \notag \\
=&\left[\left(\frac{1}{2 h^2}-\frac{1}{4h}\right) \varepsilon_{j+1}^{n}+\left(\frac{1}{\tau_{n+1}}-\frac{1-2e^{-\lambda \frac{\tau_{n+1}}{2}}}{2\Gamma(2-\alpha)(\frac{\tau_{n+1}}{2})^\alpha}-\frac{1}{h^2}\right) \varepsilon_j^{n}+\left(\frac{1}{2 h^2}+\frac{1}{4h}\right) \varepsilon_{j-1}^{n}\right] \notag\\
&+\frac{1}{\Gamma(1-\alpha)}\left[(a_{0, n}-\frac{e^{-\lambda \frac{\tau_{n+1}}{2}}}{(\frac{\tau_{n+1}}{2})^\alpha}) \varepsilon_j^n+\sum\limits_{l=1}^{n-1} \left(a_{n-l, n}+b_{n-1-l, n}\right) \varepsilon_j^l+(b_{n-1, n}+\frac{e^{-\lambda t_{n+\frac{1}{2}}}}{t_{n+\frac{1}{2}}^\alpha}) \varepsilon_j^0\right] \notag\\
&+R_j^{n+\frac{1}{2}} \label{converge1}
\end{align}
 Similar to the stability analysis, we define the grid functions as follows
 \begin{equation}\notag
\varepsilon^n(x)=\left\{\begin{array}{l}
0,  \quad 0 \leq x \leq x_{\frac{1}{2}},  \\
\varepsilon_j^n,  \quad x_{j-\frac{1}{2}} \leq x \leq x_{j+\frac{1}{2}},  \quad 1 \leq j \leq M-1,  \\
0,  \quad x_{M-\frac{1}{2}} \leq x \leq x_M . 
\end{array}\right. 
\end{equation}
\begin{equation}\notag
R^{n+\frac{1}{2}}(x)=\left\{\begin{array}{l}
0,  \quad 0 \leq x \leq x_{\frac{1}{2}},  \\
R^{n+\frac{1}{2}}_j,  \quad x_{j-\frac{1}{2}} \leq x \leq x_{j+\frac{1}{2}},  \quad 1 \leq j \leq M-1,  \\
0,  \quad x_{M-\frac{1}{2}} \leq x \leq x_M . 
\end{array}\right. 
\end{equation}
Similar to stability analysis, denote $\hat{\varepsilon}^n[k]$ and $\hat{R}^{n+\frac{1}{2}}[k]$ as the discrete fourier transform of $\varepsilon^n(x)$ and $R^{n+\frac{1}{2}}$, then we have
\begin{equation}
\begin{aligned}\label{suppose}
\hat{\varepsilon}^n[k]=\sum_{j=0}^{M-1} e^{-i \frac{2 \pi}{M} j k} \varepsilon^n_j \quad k=0, 1,  \ldots,  M-1 , \\
\hat{R}^{n+\frac{1}{2}}[k]=\sum_{j=0}^{M-1} e^{-i \frac{2 \pi}{M} j k} R^{n+\frac{1}{2}}_j \quad k=0, 1,  \ldots,  M-1 . 
\end{aligned}
\end{equation}
where \( e \) is the base of the natural logarithm,  and \( i \) is the imaginary unit.  The inverse Discrete Fourier Transform (IDFT) is given by:
\[
\varepsilon^n_j = \frac{1}{M} \sum_{k=0}^{M-1} \hat{\varepsilon}^n[k] e^{i \frac{2\pi}{M} jk}, R^{n+\frac{1}{2}}_j = \frac{1}{M} \sum_{k=0}^{M-1} \hat{R}^{n+\frac{1}{2}}[k] e^{i \frac{2\pi}{M} jk}
\]
Using the boundar conditions, it is easy to obtain $\varepsilon^n_0=\varepsilon^n_M=0$. Then we have the Parseval equality for the discrete Fourier transforms
\begin{align}
    \int_0^L\left|\varepsilon^n(x)\right|^2 d x&=h\sum_{j=0}^{M-1}\left|\varepsilon^n_j\right|^2=Mh\sum_{k=0}^{M-1}\left|\hat{\varepsilon}^n[k]\right|^2 =L\sum_{k=0}^{M-1}\left|\hat{\varepsilon}^n[k]\right|^2. \notag\\
    \int_0^L\left|R^{n+\frac{1}{2}}(x)\right|^2 d x&=h\sum_{j=0}^{M-1}\left|R^{n+\frac{1}{2}}_j\right|^2=Mh\sum_{k=0}^{M-1}\left|\hat{R}^{n+\frac{1}{2}}[k]\right|^2 =L\sum_{k=0}^{M-1}\left|\hat{R}^{n+\frac{1}{2}}[k]\right|^2. \notag
\end{align}

Introduce the following norm
$$
\left\|\varepsilon^n\right\|_{L_2}=\left(\sum_{j=1}^{M-1} h\left|\varepsilon_j^n\right|^2\right)^{1 / 2}=\left(\int_0^L\left|\varepsilon_j^n\right|^2 d x\right)^{1 / 2}
$$
$$
\left\|R^{n+\frac{1}{2}}\right\|_{L_2}=\left(\sum_{j=1}^{M-1} h\left|R^{n+\frac{1}{2}}_j\right|^2\right)^{1 / 2}=\left(\int_0^L\left|R^{n+\frac{1}{2}}_j\right|^2 d x\right)^{1 / 2}
$$
we also have
$$
\left\|\varepsilon^n\right\|_{L_2}^2=L\sum_{k=0}^{M-1}\left|\hat{\varepsilon}^n[k]\right|^2, \left\|R^{n+\frac{1}{2}}\right\|_{L_2}^2=L\sum_{k=0}^{M-1}\left|\hat{R}^{n+\frac{1}{2}}[k]\right|^2
$$
  Using the formulas in (\ref{suppose}),  We multiply both sides of equation  (\ref{converge1}) by $e^{-ijh\beta}$  where $\beta = 2\pi k/M$ and then perform summation, we obtain into
\begin{equation}
\begin{aligned}
&\sum_{j=0}^{M-1}\left[\left(-\frac{1}{2 h^2}+\frac{1}{4h}\right) \varepsilon_{j+1}^{n+1}e^{-ijh\beta}+\left(\frac{1}{\tau_{n+1}}+\frac{1}{2\Gamma(2-\alpha)(\frac{\tau_{n+1}}{2})^\alpha}+\frac{1}{h^2}\right) \varepsilon_{j}^{n+1}e^{-ijh\beta}\right. \\
&\left. +\left(-\frac{1}{2 h^2}-\frac{1}{4h}\right) \varepsilon_{j-1}^{n+1}e^{-ijh\beta}\right] \\
&=\sum_{j=0}^{M-1}\left[\left(\frac{1}{2 h^2}-\frac{1}{4h}\right) \varepsilon_{j+1}^{n}e^{-ijh\beta}+\left(\frac{1}{\tau_{n+1}}-\frac{1-2e^{-\lambda \frac{\tau_{n+1}}{2}}}{2\Gamma(2-\alpha)(\frac{\tau_{n+1}}{2})^\alpha}-\frac{1}{h^2}\right) \varepsilon_{j}^{n}e^{-ijh\beta}\right. \\
&\left. +\left(\frac{1}{2 h^2}+\frac{1}{4h}\right) \varepsilon_{j-1}^{n}e^{-ijh\beta}\right] \\
&+\sum_{j=0}^{M-1}\frac{1}{\Gamma(1-\alpha)}\left[(a_{0, n}-\frac{e^{-\lambda \frac{\tau_{n+1}}{2}}}{(\frac{\tau_{n+1}}{2})^\alpha}) \varepsilon_{j}^{n}e^{-ijh\beta}+\sum\limits_{l=1}^{n-1} \left(a_{n-l, n}+b_{n-1-l, n}\right) \varepsilon_{j}^{l}e^{-ijh\beta}\right. \\
&\left. +(b_{n-1, n}+\frac{e^{-\lambda t_{n+\frac{1}{2}}}}{t_{n+\frac{1}{2}}^\alpha}) \varepsilon_{j}^{0}e^{-ijh\beta}\right] +\sum_{j=0}^{M-1}R^{n+\frac{1}{2}} e^{-i j h \beta}
\end{aligned}
\end{equation}
After simplifications,  we have
\begin{align}
     &\left[-\frac{1}{h^2}cos(h\beta)+\frac{1}{2h}isin(h\beta)+\left(\frac{1}{\tau_{n+1}}+\frac{1}{2\Gamma(2-\alpha)(\frac{\tau_{n+1}}{2})^\alpha}+\frac{1}{h^2}\right) \right]\hat{\varepsilon}^{n+1}[k]\notag \\
    &=\left[\frac{1}{h^2}cos(h\beta)-\frac{1}{2h}isin(h\beta)+\left(\frac{1}{\tau_{n+1}}-\frac{1-2e^{-\lambda \frac{\tau_{n+1}}{2}}}{2\Gamma(2-\alpha)(\frac{\tau_{n+1}}{2})^\alpha}-\frac{1}{h^2}\right) -\frac{e^{-\lambda \frac{\tau_{n+1}}{2}}}{\Gamma(1-\alpha)(\frac{\tau_{n+1}}{2})^\alpha}\right]\hat{\varepsilon}^{n}[k] \notag\\
    &+\frac{1}{\Gamma(1-\alpha)}\left[a_{0, n} \hat{\varepsilon}^{n}[k]+\sum\limits_{l=1}^{n-1} \left(a_{n-l, n}+b_{n-1-l, n}\right) \hat{\varepsilon}^{l}[k]+(b_{n-1, n}+\frac{e^{-\lambda t_{n+\frac{1}{2}}}}{t_{n+\frac{1}{2}}^\alpha}) \hat{\varepsilon}^{0}[k]\right]\notag\\
    &+\hat{R}^{n+\frac{1}{2}}(k) \label{conLemma}
\end{align}
\begin{lemma}\label{convergele}
    Suppose that $\hat{\varepsilon}^n [k] (n = 1,  2, \cdots,  N)$ and $\hat{R}^{n+\frac{1}{2}}[k](n = 0,  1, \cdots,  N-1)$ are solutions of (\ref{conLemma}).  Then we have 
    $$
    |\hat{\varepsilon}^n[k]|\leq (1+\tau_1)^{n^{r-2}N^{2}}|\hat{R}^{\frac{1}{2}}[k]|, n = 1,  2, \cdots,  N$$
\end{lemma}
\begin{proof}
    Notice that $\varepsilon^0_j=0, j=1, 2, \cdots, M$, so $\hat{\varepsilon}^0(k)=0$. Then for $n=0$, we get 
    \begin{align}
    |\hat{\varepsilon}^{1}[k]|&= \left|\frac{1}{-\frac{1}{h^2}cos(h\beta)+\frac{1}{2h}isin(h\beta)+\left(\frac{1}{\tau_{1}}+\frac{1}{2\Gamma(2-\alpha)(\frac{\tau_{1}}{2})^\alpha}+\frac{1}{h^2}\right)}\right||\hat{R}^{\frac{1}{2}}(k)|\notag\\
    &=\frac{1}{\sqrt{\left( \frac{1}{\tau_1}+\frac{1}{2\Gamma(2-\alpha)(\frac{\tau_{1}}{2})^\alpha}+\frac{1-cos(h\beta)}{h^2}\right)^2+\left(\frac{sin(h\beta)}{2h}\right)^2}}|\hat{R}^{\frac{1}{2}}(k)|\notag\\
    &\leq |\hat{R}^{\frac{1}{2}}(k)|\notag
\end{align}
Assume that $|\hat{\varepsilon}^m(k)|\leq (1+\tau_1)^{m^{r-2}N^{2}}|\hat{R}^{\frac{1}{2}}(k)|, m=1, 2, \cdots, n$. Then using the equation (\ref{conLemma}), similar with Lemma \ref{dnd0}, we get
\begin{align}
    &\left|-\frac{1}{h^2}cos(h\beta)+\frac{1}{2h}isin(h\beta)+\left(\frac{1}{\tau_{n+1}}+\frac{1}{2\Gamma(2-\alpha)(\frac{\tau_{n+1}}{2})^\alpha}+\frac{1}{h^2}\right) \right||\hat{\varepsilon}^{n+1}(k)|\notag \\
    &\leq (1+\tau_1)^{n^{r-2}N^{2}}\left(\left|\frac{1}{h^2}cos(h\beta)-\frac{1}{2h}isin(h\beta)+\left(\frac{1}{\tau_{n+1}}-\frac{1-2e^{-\lambda \frac{\tau_{n+1}}{2}}}{2\Gamma(2-\alpha)(\frac{\tau_{n+1}}{2})^\alpha}-\frac{1}{h^2}\right) -\right. \right. \notag\\
    &\left. \left. \frac{e^{-\lambda \frac{\tau_{n+1}}{2}}}{\Gamma(1-\alpha)(\frac{\tau_{n+1}}{2})^\alpha}\right|+\frac{e^{-\lambda \frac{\tau_{n+1}}{2}}}{\Gamma(1-\alpha)(\frac{\tau_{n+1}}{2})^\alpha}\right) \left|\hat{R}^{\frac{1}{2}}(k)\right|+\left|\hat{R}^{n+\frac{1}{2}}(k)\right| \notag\\
    &\leq (1+\tau_1)^{n^{r-2}N^{2}}\left|-\frac{1}{h^2}cos(h\beta)+\frac{1}{2h}isin(h\beta)+\left(\frac{1}{\tau_{n+1}}+\frac{1}{2\Gamma(2-\alpha)(\frac{\tau_{n+1}}{2})^\alpha}+\frac{1}{h^2}\right) \right|\left|\hat{R}^{\frac{1}{2}}(k)\right|\notag\\
    &+\left|\hat{R}^{n+\frac{1}{2}}(k)\right| \label{xi}
\end{align}
Notice that $\left|R_j^{n+\frac{1}{2}}\right| \leq K[(n+1)^{-(2-\alpha)}+h^2+\varepsilon]$ ,  $\left|R_j^{\frac{1}{2}}\right| \leq K[N^{-2}+h^2+\varepsilon]$ by (\ref{R_half}) and $\hat{R}^{n+\frac{1}{2}}[k]=\sum_{j=0}^{M-1} e^{-i \frac{2 \pi}{M} j k} R^{n+\frac{1}{2}}_j $, we have
$$
\left|\hat{R}^{n+\frac{1}{2}}(k)\right|\leq \frac{N^2}{(n+1)^{2-\alpha}}\left|\hat{R}^{\frac{1}{2}}(k)\right|
$$
Substitute the above equation into (\ref{xi}) and we have
\begin{align}
    |\hat{\varepsilon}^{n+1}(k)|&\leq (1+\tau_1)^{n^{r-2}N^{2}}|\eta^{\frac{1}{2}}(k)|\notag\\
    &+\left|\frac{\frac{N^2}{(n+1)^{2-\alpha}}}{-\frac{1}{h^2}cos(h\beta)+\frac{1}{2h}isin(h\beta)+\left(\frac{1}{\tau_{1}}+\frac{1}{2\Gamma(2-\alpha)(\frac{\tau_{1}}{2})^\alpha}+\frac{1}{h^2}\right)}\right||\hat{R}^{\frac{1}{2}}(k)|\notag\\
    &=(1+\tau_1)^{n^{r-2}N^{2}}|\hat{R}^{\frac{1}{2}}(k)|+\frac{\frac{N^2}{(n+1)^{2-\alpha}}}{\sqrt{\left( \frac{1}{\tau_1}+\frac{1}{2\Gamma(2-\alpha)(\frac{\tau_{1}}{2})^\alpha}+\frac{1-cos(h\beta)}{h^2}\right)^2+\left(\frac{sin(h\beta)}{2h}\right)^2}}|\hat{R}^{\frac{1}{2}}(k)|\notag\\
    &\leq (1+\tau_1)^{n^{r-2}N^{2}}|\hat{R}^{\frac{1}{2}}(k)|+\tau_1 \frac{N^2}{(n+1)^{2-\alpha}}|\hat{R}^{\frac{1}{2}}(k)|\notag\\
    &\leq \left[(1+\tau_1)^{n^{r-2}N^{2}}+(1+\tau_1)^{\frac{N^2}{(n+1)^{2-\alpha}}}-1\right]|\hat{R}^{\frac{1}{2}}(k)|\notag\\
    &\leq(1+\tau_1)^{(n+1)^{r-2}N^{2}}|\hat{R}^{\frac{1}{2}}(k)|\notag\\
    &- \left[(1+\tau_1)^{(n+1)^{r-2}N^{2}}-(1+\tau_1)^{n^{r-2}N^{2}}-(1+\tau_1)^{\frac{N^2}{(n+1)^{2-\alpha}}}+1\right]|\hat{R}^{\frac{1}{2}}(k)|\notag
\end{align}
Denote $g = (1+\tau_1)^{\frac{N^2}{(n+1)^{2-\alpha}}}\geq 1$ and notice that $r\geq 3$,  we have
\begin{align}
    &(1+\tau_1)^{(n+1)^{r-2}N^{2}}-(1+\tau_1)^{n^{r-2}N^{2}}-(1+\tau_1)^{\frac{N^2}{(n+1)^{2-\alpha}}}+1\notag\\
    &= g^{(n+1)^{r-\alpha}}-g^{n^{r-2}(n+1)^{2-\alpha}}-g+1\notag\\
    &=g(g^{(n+1)^{r-\alpha}-1}-1)-(g^{n^{r-2}(n+1)^{2-\alpha}}-1)\notag\\
    &\geq g^{(n+1)^{r-\alpha}-1}-g^{n^{r-2}(n+1)^{2-\alpha}}\notag\\
    &= g^{n^{r-2}(n+1)^{2-\alpha}(1+\frac{1}{n})^{r-2}-1}-g^{n^{r-2}(n+1)^{2-\alpha}}\notag\\
    &\geq g^{n^{r-2}(n+1)^{2-\alpha}(1+\frac{r-2}{n})-1}-g^{n^{r-2}(n+1)^{2-\alpha}}\notag\\
    &\geq g^{n^{r-2}(n+1)^{2-\alpha}+(r-2)n^{r-3}(n+1)^{2-\alpha}-1}-g^{n^{r-2}(n+1)^{2-\alpha}}\notag\\
    &\geq 0 \notag
\end{align}
So we get $|\hat{\varepsilon}^{n+1}(k)|\leq (1+\tau_1)^{(n+1)^{r-2}N^{2}}\left|\hat{R}^{\frac{1}{2}}(k)\right|$. By using mathematical induction,  we get $|\hat{\varepsilon}^n(k)|\leq (1+\tau_1)^{n^{r-2}N^{2}}|\hat{R}^{\frac{1}{2}}(k)|, n=1, 2, \cdots, N$. So we complete the proof. 
\end{proof}
Now we can give the proof of theorem \ref{converge}. 
    \begin{proof}
    Using Lemma \ref{convergele} ,  we obtain
\begin{align}
    |\hat{\varepsilon}^n(k)|&\leq (1+\tau_1)^{n^{r-2}N^{2}}|\hat{R}^{\frac{1}{2}}(k)|\notag\\
    &\leq (1+\frac{T}{N^r})^{N^{r-2}N^{2}}|\hat{R}^{\frac{1}{2}}(k)|\notag\\
    &=(1+\frac{T}{N^r})^{N^r}|\hat{R}^{\frac{1}{2}}(k)|\notag\\
    &\leq \lim_{N\to \infty}(1+\frac{T}{N^r})^{N^r}|\hat{R}^{\frac{1}{2}}(k)|\notag\\
    &=e^T|\hat{R}^{\frac{1}{2}}(k)|\notag
\end{align}
    So using the Parseval equality and let $C=e^T$, we obtain
    \begin{equation}
\begin{aligned}
\left\|u^n-U^n\right\|_{L_2}^2 & =\left\|\varepsilon^n\right\|_{L_2}^2=L\sum_{k=0}^{M-1}\left|\hat{\varepsilon}^n[k]\right|^2\\
&\leq C^2L \sum_{k=0}^{M-1}\left|\hat{R}^{\frac{1}{2}}[k]\right|^2=C^2\left\|R^{\frac{1}{2}}\right\|_{L_2}^2\leq C^2KL(N^{-2}+h^2+\varepsilon)^2. 
\end{aligned}
\end{equation}
which completes the proof. 
\end{proof}

\section{Numerical Experiments}
\subsection{Example 1}

Let us start with the numerical differentiation formulas for time Caputo tempered fractional derivatives ${}^C_0D_t^{\alpha, \lambda}u(t)$. Let $T = 2$ and $u(t) = t^\delta$.  Use the formulas of numerical differentiation to find ${}^C_0D_t^{\alpha, \lambda}u(t_{n+\frac{1}{2}})$. Considering the difference
$$
E_{max}(N)=\max_{0\leq n \leq N-1}|{}^C_0D_t^{\alpha, \lambda}u(t_{n+\frac{1}{2}})-{}^{FC}_0D_t^{\alpha, \lambda}u(t_{n+\frac{1}{2}})|
$$
where $N$ is a temporal step number and $t_n = (n/N)^r T, t_{n+\frac{1}{2}}=\frac{t_n+t_{n+1}}{2}$,  we define the truncation error
order by
$$
order = log_2\frac{E_{max}(N/2)}{E_{max}(N)}
$$
\begin{table}[h!]
\centering
\begin{tabular}{|c|c|c|c|c|c|c|}
\hline
$N$ & \multicolumn{2}{|c|}{$\alpha=0. 1$} & \multicolumn{2}{|c|}{$\alpha=0. 3$} & \multicolumn{2}{|c|}{$\alpha=0. 5$} \\ \hline
 & $E_{max}(N)$ & order & $E_{max}(N)$ & order & $E_{max}(N)$ & order \\ \hline
80  & 8. 1838e-05 & -   & 5. 8934e-04 & -   & 3. 0071e-03 & -   \\ \hline
160  & 2. 1028e-05 & 1. 9605 & 1. 8795e-04 &  1. 6488&1. 0835e-03  & 1. 4727 \\ \hline
320 & 5. 4093e-06 & 1. 9588 & 5. 9220e-05 & 1. 6662 & 3. 8718e-04 & 1. 4846 \\ \hline
640 & 1. 3932e-06 & 1. 9570 & 1. 8528e-05 & 1. 6764 & 0. 0019471 & 1. 4910 \\ \hline
\end{tabular}
\caption{Error table for $u(t) = t^\delta, r=1. 5,  \delta=1. 5$}
\label{table:errors}
\end{table}

\subsection{Example 2}

Next,  we demonstrate the convergence order and computational complexity of the Fast
Scheme. 

Consider the following problem
\begin{equation}\label{example}
    \begin{aligned}
        \left\{\begin{array}{l}
\frac{\partial u(x,  t)}{\partial t}+{ }_0^C D_t^{\alpha,  \lambda} u(x,  t)=\frac{\partial^2 u(x,  t)}{\partial x^2}- \frac{\partial u(x,  t)}{\partial x}+f(x,  t),  \quad(x,  t) \in[0, 1] \times[0, 2],  \\
u(x,  0)=x^2(1-x)^2,  \quad x \in[0, 1],  \\
u(0,  t)=u(1,  t)=0,  \quad t \in[0, 2], 
\end{array}\right. 
    \end{aligned}
\end{equation}
where $0<\alpha<1,  \lambda=1,  \delta=1. 8$ and
$$
\begin{aligned}
f(x,  t)= & \left(-\lambda\left(t^\delta+1\right)+\delta t^{\delta-1}+\frac{\Gamma(\delta+1)}{\Gamma(\delta-\alpha+1)} t^{\delta-\alpha}\right) e^{-\lambda t} x^2(1-x)^2 \\
& -\left(\left(12 x^2-12 x+2\right)-\left(4 x^3-6 x^2+2 x\right)\right)\left(t^\delta+1\right) e^{-\lambda t} . 
\end{aligned}
$$
The exact solution of (\ref{example}) is $u(x,  t)=e^{-\lambda t}\left(t^\delta+1\right) x^2(1-x)^2$. Denote $U_i^n=U(x_i, t_n)(0\leq i\leq M, 0\leq n\leq N)$ as the numerical solutions and considering the difference
$$
e_{max}(M, N) = \max_{0\leq n\leq N}\sqrt{h\sum_{i=0}^M(u(x_i, t_n)-U_i^n)^2}, 
$$
$$order = \frac{log(\frac{e_{max}(M, N)}{e_{max}(M, 2N)})}{log(\frac{\tau_N}{\tau_{2N}})}$$
\begin{table}[h!]
\centering
\begin{tabular}{|c|c|c|c|c|c|c|}
\hline
$N$ & \multicolumn{2}{|c|}{$\alpha=0. 1$} & \multicolumn{2}{|c|}{$\alpha=0. 3$} & \multicolumn{2}{|c|}{$\alpha=0. 5$} \\ \hline
 & $e_{max}(M, N)$ & order & $e_{max}(M, N)$ & order & $e_{max}(M, N)$ & order \\ \hline
10  & 1. 3449e-03 & -   & 1. 3568e-03 & -   & 1. 3802e-03 & -  \\ \hline
20  & 3. 3323e-04 &2. 1736 & 3. 3372e-04 & 2. 1688
 &3. 4626e-04  & 2. 1542 \\ \hline
40 & 8. 3166e-05 & 2. 0783 & 8. 4490e-05 & 2. 0724
 & 8. 7981e-05 & 2. 0515 \\ \hline
80 & 2. 0814e-05 & 2. 0353 & 2. 1233e-05 &2. 0292
  & 2. 2555e-05 & 2. 0000 \\ \hline
  160 & 5. 2056e-06 & 2. 0177 & 5. 3435e-06 &2. 0087
  & 5. 8399e-06 & 1. 9672 \\ \hline
\end{tabular}
\caption{Error table where $r=3, \delta=1. 8$}
\label{table:errors}
\end{table}
\section{Conclusion}
In this paper,  we develop a novel fast finite‐difference scheme for the tempered time‐fractional derivative exhibiting weak regularity at the initial time.   The scheme employs piecewise linear interpolation over each small temporal subinterval and yields a fully discrete approximation for the advection–dispersion equation.   To raise the overall accuracy,  temporal differencing is carried out at the half‐time levels,  while spatial derivatives are approximated by centered differences.   In contrast to earlier work,  our method achieves second‐order convergence in both time and space,  which is a breakthrough in the temporal discretization of tempered fractional models.   

In addition, concerning possible optimizations and extensions of the proposed approach, we would like to refer to \cite{zhao2021}, where an all-at-once system is established. This formulation is well suited for parallel computing and provides an alternative strategy for fast numerical solution. Regarding \cite{Alikhanov2024_1,Alikhanov2024_2,Alikhanov2024_3,Vabishchevich2022,Vabishchevich2023}, they consider general problems involving difference kernels, which may offer a potential avenue for extending the current model to more general kernel formulations.
\vspace{18pt}

\noindent {\LARGE \bf Declarations}

\vspace{12pt}

\noindent{\bf Conflict of interest}\;\;
 The authors declare that they have no conflict of interest. 


\vspace{18pt}


\begin{thebibliography}{1}
\bibitem{Alikhanov2024_1}
A. A. Alikhanov, M. S. Asl and C. Huang, Stability analysis of a second-order difference scheme for the time-fractional mixed sub-diffusion and diffusion-wave equation, Fractional Calculus and Applied Analysis, Vol. 27 Issue 1 Pages 102-123 (2024)

\bibitem{Alikhanov2024_2}
A. A. Alikhanov, M. S. Asl, C. Huang and A. Khibiev, A second-order difference scheme for the nonlinear time-fractional diffusion-wave equation with generalized memory kernel in the presence of time delay, Journal of Computational and Applied Mathematics, Vol. 438 Pages 115515 (2024)

\bibitem{Alikhanov2024_3}
A. A. Alikhanov, M. S. Asl and D. Li, A novel explicit fast numerical scheme for the Cauchy problem for integro-differential equations with a difference kernel and its application, Computers \& Mathematics with Applications, Vol. 175 Pages 330-344 (2024)

\bibitem{beylkin2005}
G.  Beylkin and L.  Monzón,  On approximation of functions by exponential sums,  Applied and Computational Harmonic Analysis,  Vol.  19 Issue 1,  Pages 17-48 (2005)

\bibitem{cao2020}
J.  Cao,  A.  Xiao and W.  Bu,  Finite Difference/Finite Element Method for Tempered Time Fractional Advection–Dispersion Equation with Fast Evaluation of Caputo Derivative,  Journal of Scientific Computing,  Vol.  83 (2020)

\bibitem{cao2023}
J.  Cao,  A.  Xiao and W.  Bu,  A fast Alikhanov algorithm with general nonuniform time steps for a two-dimensional distributed-order time–space fractional advection–dispersion equation,  Numerical Methods for Partial Differential Equations,  Vol.  39 Issue 4,  Pages 2885-2908 (2023)

\bibitem{carr2002}
P.  Carr,  H.  Geman,  D.  Madan and M.  Yor,  The Fine Structure of Asset Returns: An Empirical Investigation,  The Journal of Business,  Vol.  75 (2002)

\bibitem{carr2003}
P.  Carr,  H.  Geman,  D.  Madan and M.  Yor,  Stochastic Volatility for Lévy Process,  Mathematical Finance,  Vol.  13,  Pages 345-382 (2003)

\bibitem{cartea2007}
Á.  Cartea and D.  del-Castillo-Negrete,  Fluid limit of the continuous-time random walk with general L\'evy jump distribution functions,  Physical Review E,  Vol.  76 Issue 4,  Pages 041105 (2007)

\bibitem{diethelm2010}
K.  Diethelm,  The Analysis of Fractional Differential Equations: An Application-Oriented Exposition Using Differential Operators of Caputo Type, Springer Verlag,  Vol.  2004 (2010)

\bibitem{feng2022}
L.  Feng,  F.  Liu,  V.  V.  Anh and S.  Qin,  Analytical and numerical investigation on the tempered time-fractional operator with application to the Bloch equation and the two-layered problem,  Nonlinear Dynamics,  Vol.  109 Issue 3,  Pages 2041-2061 (2022)

\bibitem{fenwick2024}
J.  Fenwick,  F.  Liu and L.  Feng,  New insight into the nano-fluid flow in a channel with tempered fractional operators,  Nanotechnology,  Vol.  35 Issue 8,  Pages 085403 (2024)

\bibitem{gracia2018}
J.  L.  Gracia,  E.  O'Riordan and M.  Stynes,  A Fitted Scheme for a Caputo Initial-Boundary Value Problem,  Journal of Scientific Computing,  Vol.  76 (2018) 

\bibitem{gregory1969}
R. T.  Gregory and D. L.  Karney,  A Collection of Matrices for Testing Computational Algorithms.  Wiley-Inter-science,  John Wiley \& Sons,  New York,  London ,  Sydney,  Toronto.  (1969)

\bibitem{koponen1995}
I.  Koponen,  Analytic approach to the problem of convergence of truncated L\'evy flights towards the Gaussian stochastic process,  Physical Review E ,  Vol.  52 Issue 1 Pages 1197-1199 (1995)

\bibitem{krzyzanowski2024}
G.  Krzyzanowski and M.  Magdziarz,  A tempered subdiffusive Black–Scholes model,  Fractional Calculus and Applied Analysis,  Vol.  27 (2024)

\bibitem{dwivedi2024}
H.  Kumar Dwivedi and Rajeev,  A novel fast tempered algorithm with high-accuracy scheme for 2D tempered fractional reaction-advection-subdiffusion equation,  Computers and Mathematics with Applications,  Pages 371-397 (2024)

\bibitem{dwivedi2025}
H.  Kumar Dwivedi and Rajeev,  A novel fast second order approach with high-order compact difference scheme and its analysis for the tempered fractional burgers equation,  Mathematics and Computers in Simulation,  Vol.  13,  Pages 1-25 (2025)

\bibitem{mantegna1994}
R.  N.  Mantegna and H.  E.  Stanley,  Stochastic Process with Ultraslow Convergence to a Gaussian: The Truncated L\'evy Flight,  Physical Review Letters,  Vol.  73 Issue 22,  Pages 2946-2949 (1994)

\bibitem{meerschaert2008}
M.  Meerschaert,  Y.  Zhang and B.  Baeumer,  Tempered anomalous diffusion in heterogeneous systems,  Geophysical Research Letters,  Vol.  35,  Pages L17403 (2008)

\bibitem{morgado2019}
L.  Morgado and L.  Morgado,  Modelling Time-of-Flight Transient Currents with Time-Fractional Diffusion Equations, Progress Fract.  Differ.  Appl.  (2019)

\bibitem{Vabishchevich2022}
P.N. Vabishchevich, Numerical solution of the Cauchy problem for Volterra integrodifferential equations with difference kernels, Applied Numerical Mathematics, Vol. 174 Pages 177-190 (2022)

\bibitem{Vabishchevich2023}
P.N. Vabishchevich, Numerical-analytical methods for solving the Cauchy problem for evolutionary equations with memory, Lobachevskii Journal of Mathematics, Vol. 44 Issue 10 Pages 4195-4204 (2023)

\bibitem{qiao2022}
L.  Qiao,  J.  Guo and W.  Qiu,  Fast BDF2 ADI methods for the multi-dimensional tempered fractional integrodifferential equation of parabolic type,  Computers \& Mathematics with Applications,  Vol.  123,  Pages 89-104 (2022)

\bibitem{safari2022}
Z.  Safari,  G.  B.  Loghmani and M.  Ahmadinia,  Convergence Analysis of a LDG Method for Time–Space Tempered Fractional Diffusion Equations with Weakly Singular Solutions,  Journal of Scientific Computing,  Vol.  91 Issue 2,  Pages 68 (2022)

\bibitem{sun2006}
Z. -z.  Sun,  A fully discrete scheme for a diffusion-wave system,  Applied Numerical Mathematics - APPL NUMER MATH,  Vol.  56,  Pages 193-209 (2006)

\bibitem{wang2022}
H.  Wang and C.  Li,  Fast difference scheme for a tempered fractional Burgers equation in porous media,  Applied Mathematics Letters,  Vol.  132,  Pages 108143 (2022)

\bibitem{xia2013}
X.  I.  A.  Yuan,  W.  U.  Jichun and Z.  Yong,  Tempered time-fractional advection-dispersion equation for modeling non-Fickian transport,  Advances in Water Science,  Vol.  24 Issue 3,  Pages 349-357 (2013)

\bibitem{zhang2011}
Y.  Zhang and M.  Meerschaert,  Gaussian setting time for solute transport in fluvial systems,  Water Resour.  Res,  Vol.  47 (2011)

\bibitem{zhang2012}
Y.  Zhang,  M.  Meerschaert and A.  Packman,  Linking fluvial bed sediment transport across scales,  Geophysical Research Letters,  Vol.  39,  Pages 20404 (2012)

\bibitem{zhang2014}
Y.  Zhang,  L.  Chen,  D.  Reeves and H.  Sun,  A fractional-order tempered-stable continuity model to capture surface water runoff,  Journal of Vibration and Control,  Vol.  22,  Pages 1993-2003 (2014)

\bibitem{zhao2020}
L.  Zhao,  C.  Li and F.  Zhao,  Efficient Diference Schemes for the Caputo-Tempered Fractional Difusion Equations Based on Polynomial Interpolation,  Communications on Applied Mathematics and Computation (2020)

\bibitem{zhao2023}
T.  Zhao,  Efficient Spectral Collocation Method for Tempered Fractional Differential Equations,  Fractal and Fractional,  Vol.  7 Issue 3,  Pages 277 (2023)

\bibitem{zhao2024}
T.  Zhao and L.  Zhao,  Efficient Jacobian Spectral Collocation Method for Spatio-Dependent Temporal Tempered Fractional Feynman-Kac Equation,  Communications on Applied Mathematics and Computation (2024)

\bibitem{zhao2021}
Y.-L. Zhao, X.-M. Gu, M. Li and H.-Y. Jian, Preconditioners for all-at-once system from the fractional mobile/immobile advection–diffusion model, Journal of Applied Mathematics and Computing, Vol. 65 Issue 1 Pages 669-691 (2021)

\bibitem{zhou2023}
J.  Zhou,  X. -M.  Gu,  Y. -L.  Zhao and H.  Li,  A fast compact difference scheme with unequal time-steps for the tempered time-fractional Black–Scholes model,  International Journal of Computer Mathematics,  Vol.  101,  No.  9-10,  Pages 989- 1011 (2024)

\end{thebibliography}


\end{document}


